\pdfmapfile{+lm.map}
\pdfmapfile{+cm.map}
\pdfmapfile{+symbols.map}
\documentclass[11pt,a4paper]{article}
\usepackage[margin=24mm]{geometry}
\usepackage[T1]{fontenc}
\usepackage{lmodern,amsmath,amssymb,amsthm,mathtools,booktabs,longtable,array,microtype}
\usepackage{xcolor,enumitem,fancyvrb,fvextra}
\usepackage[hidelinks,breaklinks]{hyperref}
\usepackage{bookmark}
\newcommand{\Z}{\mathbb Z}\newcommand{\N}{\mathbb N}
\newcommand{\Con}{\operatorname{Con}}
\newcommand{\tightlist}{\setlength{\itemsep}{0pt}\setlength{\parskip}{0pt}}
\newtheorem{theorem}{Theorem}[section]
\newtheorem{lemma}[theorem]{Lemma}
\theoremstyle{remark}\newtheorem{remark}[theorem]{Remark}
\setlength{\emergencystretch}{3em}
\setlist{itemsep=3pt,topsep=5pt}
\DefineVerbatimEnvironment{verbatim}{Verbatim}{fontsize=\small,breaklines=true,breakanywhere=true,breaksymbolleft={},xleftmargin=1em}
\hypersetup{pdftitle={Diophantine consistency equations with unified binomials and even Pell indices},pdfauthor={Research contribution prepared for Vladimir Reshetnikov}}
\title{Diophantine consistency equations with unified binomials and even Pell indices\\[8pt]\large Round 15 research report and reproducible artifacts}
\author{Prepared for Vladimir Reshetnikov and the ProveIt repository}
\date{24 October 2026}
\begin{document}
\maketitle
\begin{abstract}
This iteration improves two fixed-polynomial constructions whose unsolvability, under explicitly retained ordinary calibration theorems, is equivalent in ZFC to $\Con(\mathrm{ZFC})$. The compact branch gives eleven unrestricted integer unknowns at exact degree $416$ with $125$ live binary arithmetic operations. Its degree/cost companions are $456/124$, $464/119$, and $488/118$. It composes coprime guard fusion, powered divisibility saturation, supported-modulus Pell interpolation, three-adic radical separation, and a cancellation-aware paired norm. The direct branch identifies the factorial denominator and reciprocal CRT product as one binomial. A new signed compilation uses selectively doubled Pell indices and remainder residues to orient powers; its principal quartic has $442$ unrestricted integer unknowns, $4652$ collected monomials, and coefficient height $128446$. A companion natural-domain quartic has $329$ nonnegative unknowns. These are measured improvements against pinned parents, not claims of global minimality. All new unbounded proofs, exact equations, complete coefficient payloads, source pins, licenses, and independent finite checks accompany this report. No new end-to-end Lean or Rocq endpoint or proof of consistency is claimed.
\end{abstract}
\clearpage
\tableofcontents
\clearpage
\section{Results and their precise interpretation}
Throughout, $\N=\{0,1,2,\ldots\}$, whereas a positive integer is strictly greater than zero. All $\Z$ endpoints quantify every listed unknown over unrestricted integers. A named register is only a polynomial abbreviation. Coefficient height means the maximum absolute coefficient after complete collection; the number of residual rows is not an arithmetic-operation count.

\begin{center}\small
\begin{tabular}{@{}lrrrrr@{}}\toprule
Endpoint and domain&Variables&Degree&Rows&Terms&Height\\\midrule
Split carrier centered, $\Z$&11&416&---&not expanded&not measured\\
Split carrier direct, $\Z$&11&456&---&not expanded&not measured\\
Centered norm, $\Z$&11&464&---&not expanded&not measured\\
Direct norm, $\Z$&11&488&---&not expanded&not measured\\
Direct certificate, $\Z$&442&4&290&4652&128446\\
Derived labels, $\N$&329&4&253&12210&86844780\\\bottomrule
\end{tabular}\end{center}
The split-carrier pair costs $125$ and $124$ operations, respectively, and trades a lower degree for a slightly higher cost. The four selected degree/cost points are $(416,125)$, $(456,124)$, $(464,119)$, and $(488,118)$. The lower-cost centered norm costs $119$ live binary operations: $74$ multiplications, $24$ additions, and $21$ subtractions. Its direct-coordinate companion costs one addition less. Squares and multiplication by scalar constants are charged; register reuse and fixed leaves are free. Neither polynomial has an unbound coefficient port. The six source34 payloads contain $4\,927\,026$ significant bits in total and $615\,883$ minimal raw bytes. Their large size is not hidden by the arithmetic-operation metric.

Against the incoming eleven-integer $688/126$ and $728/122$ degree/cost points, the selected norm pair improves both degree and cost. Against the delivered round-14 $519$-integer quartic, the principal direct equation reduces the variables from $519$ to $442$, residuals from $321$ to $290$, monomials from $6253$ to $4652$, and height from $411060$ to $128446$. The $329$-natural equation improves the pinned $349$-natural parent in all four corresponding measures. Cross-domain comparisons are not used as dominance claims.

\subsection{What the conclusion says}
For each displayed endpoint $P$ on its stated domain $D$, the conclusion is
\[
 \mathrm{ZFC}\vdash
 \left[\forall x\in D^n\;P(x)\ne0\right]
 \longleftrightarrow\Con(\mathrm{ZFC}).
\]
The ordinary proof goes through
\[
 (\exists x\in D^n\;P(x)=0)
 \longleftrightarrow
 \text{there is a finite syntactic ZFC refutation}.
\]
No consistency assumption or proof of either side is asserted. Finite reconstruction proves statements about the supplied literal polynomials and certificates, not their unsolvability. The unbounded arithmetic reasoning is an ordinary proof, with inherited formalizability and logical calibration explicitly retained. There is no newly completed end-to-end Lean, Rocq, or encoded object-ZFC proof object in this release.

\subsection{Source pins and calibration boundaries}
The public comparison snapshot is ProveIt commit
\begin{center}\small\texttt{45b399ec51dfbecb891dd51d8a9e4b6b7c361ac6}.\end{center}
It was committed at 22:28:29 UTC on 24 October 2026 and adds eight incoming ZIP files. Addition to that directory is not integration or acceptance into active source. The active public equation still has ten strictly positive unknowns, exact degree $360$, $98$ live operations, and $295388$ raw coefficient bytes; it uses a different coherent tuple. The incoming positive frontier and unrestricted-integer frontier are distinguished below.

The original coefficient recipe, imported round-14 comparison, and underlying direct C222 calibration retain the following respective pins:
\begin{center}\small\begin{tabular}{l}\texttt{2605669a532342f153b13eb6b4a0af4b54f5fa37}\\
\texttt{84e16e43092614bae262ca59651cb9d77e31a255}\\
\texttt{08200eb10fde2db408c69dffd840d6920a569b1f}\end{tabular}\end{center} A new comparison pin does not rewrite the provenance of inherited code or proofs.

The direct logical interfaces are named
\texttt{KEndpoint.c222\_iff\_zf\_refutation} and
\texttt{ConstructibleChoice.zf\_refutation\_iff\_zfc\_refutation}.
The direct route does not simulate K215 or instantiate a universal-polynomial black box. Its simultaneous Pell power graph, source-rule syntax, Horn-flow interface, and exact natural contraction are inherited. The scalar route retains the source34/native/compiler/refutation calibration, adjacent-coefficient descent, signed folding, pre-projection rank bounds, and canonical outer witnesses. Complete retained proof texts and license notices are bundled. The source34 six-port tuple must not be replaced by the compact44/BFS tuple from the positive-variable lane.

\subsection{A map of the new argument}
The scalar branch starts with all eleven inputs signed. A nonzero carrier and a disk budget establish the hypotheses needed for three-adic radical separation. Coprimality splits the fused guard square. Signed Pell folding recovers ranks before fifth-power divisibility is used, and a strict separation bound restores the exact main projection. A single strong Pell pair supports every required auxiliary residue; completeness never changes it after choosing that residue. A paired norm identity, exact centering cancellation, and absorption of a constant guard weight produce the measured degree and operation counts.

The direct branch first treats the proposed binomial output as uninterpreted. Raw ceiling rows and an independently valid power give $V<B^2$, so the pure radix is large enough to recover that binomial without circularity. A padded-horizon fixed-point proof then forces $F=t!$. The same binomial is the product of the reciprocal CRT moduli. Carry-free digits and valuation bounds decode every row. The signed compiler keeps two seed indices ordinary, doubles thirteen others, and enforces residue one through eleven existing remainders. The sign bootstrap precedes orientation. Completeness rebuilds private Pell witnesses rather than pretending the changed Pell certificate is a coordinate bijection.

\subsection{Literal conventions}
The scalar files are finite directed acyclic graphs over binary $+,-,\times$, integer literals, eleven variable leaves, and six authenticated fixed-coefficient leaves. Their output node defines the polynomial. For direct endpoints a term $[c,[i_1,\ldots,i_k]]$ denotes $c\prod_j x_{i_j}$ with zero-based variable indices; the empty index list denotes a constant. The equation is the sum of squares of \emph{every} listed residual. Full collected coefficient lists are supplied.

The complete proof sections below use readable register notation when a long expression mirrors the executable artifact. In those displays, a caret denotes exponentiation; variable exponents occur only in proof explanations, never as primitives in a final polynomial. An asterisk denotes ordinary multiplication, and \texttt{sqrt} occurs only in proofs of polynomial identities or witness recovery, never as a primitive in a final equation. References to a source proof mark an inherited boundary; all new unbounded arithmetic steps are reproduced in this report rather than deferred to a Markdown file.

\section{The fixed scalar equations in mathematical notation}
This section gives a typeset reference for the literal register displays in the complete proofs that follow. The independent unknowns are
\[
 (j,C,z,h,r,v,w,T,c,n,s)\in\Z^{11}.
\]
The fixed source34 coefficients are $\beta,K_0,M_C,M_F,A_{\rm cone},B_{\rm cone}$, with
\[
 \beta=2^{1064921}-1,\qquad A_{\rm cone}=2^{1113},\qquad B_{\rm cone}=2^{666104}.
\]
The other three integers are completely specified by the included authenticated byte payloads and printed-seed recipe. In particular $M_F$ is already shifted by $\beta$. No coefficient is an additional quantified unknown.

\subsection{The lower cost pair}
Define the following polynomial abbreviations:
\begin{align*}
 J&=6j+1,&q&=\beta J+1,&W&=12w+4,\\
 X&=qW,&Y&=(2v+1)q^3,&E_0&=XY,\\
 a&=E_0+Y,&H&=4a+3,&\Delta&=a^2+H.
\end{align*}
Then set
\begin{align*}
 F&=(K_0+W)C-z(q-1),\\
 R&=(q(q-F)-C)(q^2-1)+(M_C+qM_F)J,\\
 k&=R+1+hE_0,\qquad \eta=c-Yk,\\
 D_{\rm err}&=(A_{\rm cone}F-q)^2+(B_{\rm cone}C-q)^2,\\
 \mathcal B&=k^2(q^2-D_{\rm err})-4q^2\eta^2,\\
 D&=3\mathcal B-4,\qquad L=XY^2k,\qquad D_t=1+L(L+k),\\
 D_g&=1+DD_ts^2,\qquad P_g=D_tD_g,\\
 I&=rc^2,\qquad D_f=1+\Delta I^2,\qquad Q=\Delta^2 I^2,\qquad\delta=Q-1,\\
 V&=D_f(c-R+4IT)-c,\qquad C_y=QV^2-1,\\
 K&=X(3c-2X)+2,\qquad z_c=HT+K,\\
 A_c&=(z_c^2+3H^2)(z_c^3-19H^3),\qquad S=srH,\\
 A&=9S^2D_f,\qquad B=S^2\delta C_y,\qquad G=81S^2P_g.
\end{align*}
The centered coordinate is $Z=S(n+QV)-A_c$; the direct coordinate is $Z=Sn-A_c$. In both cases the final polynomial is
\begin{equation}\label{eq:paired}
 \mathcal P=
 \left[(Z^2+A-B-G)^2-4(AZ^2+BG)\right]^2
 -64(AZ^2)(BG).
\end{equation}
Every display is a polynomial abbreviation; powers here have fixed nonnegative integer exponents. Expanding the abbreviations gives a single fixed ordinary integer polynomial in exactly eleven unknowns.

The paired identity defining the proof interface is
\begin{align*}
 \mathcal P
 &=\left[(Z^2-A-B-G)^2-4(AB+AG+BG)\right]^2-64ABGZ^2\\
 &=\prod_{\epsilon\in\{-1,1\}^3}
   (Z-\epsilon_1\sqrt A-\epsilon_2\sqrt B-\epsilon_3\sqrt G).
\end{align*}
The last equality is a universal polynomial identity, not an assumption about radical independence or signs. The actual graph evaluates the first, polynomial-only expression. Its weighted guard uses
\[
 D_{81}=243\mathcal B-324,\quad
 D_{g,81}=81+D_{81}D_ts^2=81D_g,\quad
 P_{81}=D_tD_{g,81},\quad G=S^2P_{81},
\]
which removes one scalar multiplication without changing the polynomial.

For the centered degree calculation put $U=Sn-A_c$. The exact cancellation is
\begin{equation}\label{eq:center}
 Z^2-B=U^2+2SQVU+S^2\bigl[Q(V^2+1)-1\bigr].
\end{equation}
The lower-cost degrees are $\deg A=34$, $\deg B=122$, $\deg G=110$, and the right side of \eqref{eq:center} has degree $96$. Thus the paired inner bracket has top form $-4B_*G_*$ and the whole centered norm has top form $16B_*^2G_*^2$. One attained coefficient is
\[
 [j^{260}w^{56}v^{64}h^{12}s^{12}r^{24}c^{32}C^4]\mathcal P
 =2^{460}3^{326}\beta^{260}>0.
\]
Its exponents sum to $464$. Without centering, $B_*^4$ is the highest homogeneous form, of degree $488$. The full proofs below establish the global bounds and authentic-coefficient attainment, not just a sampled degree.

\subsection{The lower degree split carrier pair}
Keep the prefix and norm in \eqref{eq:paired}, but change $Y$ to $qv$ and replace the carrier definitions by
\begin{align*}
 M&=2q^2,& D_c&=HM,& S&=srHM^5,\\
 z_c&=D_cT+MK+H(v-q^2),&&&\\
 A_c&=(z_c^2+3D_c^2)(z_c^3-19D_c^3).&&&
\end{align*}
The ordinary proof must not assume the old cubic scale at this point. Instead it first obtains integral radicals and $S\mid A_c$, even when $c=0$ or $v=0$. Since $M^5\mid A_c$ and $M\mid D_c$, a least-valuation argument forces $M\mid z_c$. The unconditional identity $\gcd(H,M)=1$ then gives
\[
 v=q^2(2v_{\rm old}+1),\qquad
 v_{\rm old}=\frac{v-q^2}{2q^2}\in\Z,
 \qquad Y=(2v_{\rm old}+1)q^3.
\]
Only after recovering that scale may one exclude $c=0$ and apply adjacent-Pell descent. The separate $H$ component of the carrier still yields only fifth-power main projection, which is handled by the same quantitative saturation argument. Completeness satisfies both carrier components by finite CRT and the supported-modulus interpolation theorem with the strong pair held fixed.

Here $\deg A=46$, $\deg B=114$, $\deg G=94$, and \eqref{eq:center} has degree $92$. The centered norm consequently has degree $416$, with an attained coefficient
\[
 [j^{224}w^{56}v^{52}C^{16}s^{12}r^{24}c^{32}]\mathcal P
 =2^{400}3^{290}\beta^{224}>0.
\]
The direct companion has highest homogeneous form $B_*^4$ and degree $456$. Its full leading-form proof keeps both degree-five summands of the leading part of $k$; discarding either would be incorrect.

\subsection{Key direct identities}
For the natural architecture the intended meanings are
\[
 \rho=F^2,\qquad V_0=(F^2+1)^t,\qquad
 B_0=\binom{F^2+t}{t},\qquad K_0'=2R_0,\qquad
 N=B_0^{K_0'},\qquad q=N^{K_0'}.
\]
Here $R_0=2+t+x+y+F$, and $K_0'$ is the direct amplification exponent, not the fixed scalar coefficient $K_0$ above; likewise $B_0,V_0,R_0$ are local direct-proof names. The raw ceiling rows, before binomial interpretation, are
\[
 V_0+B_0=B_0F+u+1,\qquad B_0=u+v+1.
\]
They imply $V_0+v=B_0F$, then $F\ge1$, $B_0>F$, and $V_0<B_0^2$. This orders the radix bootstrap before exact binomial recovery. The fixed-point theorem forces $F=t!$ for the padded horizon $t\ge16$, after which
\[
 \prod_{i=1}^{t}\left(1+\frac{F^2}{i}\right)
 =\binom{F^2+t}{t}=B_0.
\]
The signed compiler's residue identity is
\[
 m^2=(q+1)^{2k},\quad m\equiv1\pmod q,\quad q\ge3
 \quad\Longrightarrow\quad m=(q+1)^k.
\]
It is implemented without an extra unknown by the existing remainder replacement $r=1+qu$. This is an existence-preserving restriction after power semantics, not a bijection on unrestricted coordinate spaces. The complete proof below establishes the needed order of signs, powers, and CRT decoding.

% Complete proof from scalar/PROOF.md; literal text, no Markdown math interpretation.
\section{Eleven unrestricted integers, a powered scalar carrier, and degree 464}



\subsection*{Result and scope}



The selected equation has eleven unrestricted integer unknowns, the six unchanged source34 coefficient payloads, exact total degree 464, and a literal 119-operation arithmetic circuit. Its direct companion has degree 488 and 118 operations. These improve the inspected selective-fusion parent, respectively degree 688 at 126 operations and degree 728 at 122 operations. The new reduction combines selective coprime fusion with a powered-divisibility carrier and genuinely stronger supported-modulus Pell interpolation. Switching from parity to prime-three separation permits the removal of two evenness factors. It is not merely a refactoring of the old final norm.



The conclusion, composed with the precisely retained ordinary native/\allowbreak{}compiler/\allowbreak{}refutation calibration, is that ZFC proves



\begin{verbatim}
Con(ZFC) iff the displayed fixed polynomial has no zero in Z^11.
\end{verbatim}

This construction neither proves either side nor supplies a new end-to-end Lean or internal object-ZFC certificate. It makes no global minimum or historical-priority claim. The full literal circuit and cancellation/\allowbreak{}degree audit are supplied separately. Their finite algebra checks do not replace the unbounded proofs below.



\subsection*{1. Fixed coefficients and complete polynomial}



The independent variables are exactly



\begin{verbatim}
j, C, z, h, r, v, w, T, c, n, s,
\end{verbatim}

all ranging over Z without any sign restrictions. The six fixed coefficients are exactly beta,K0,MC,MF,cone\_\allowbreak{}A,cone\_\allowbreak{}B from coefficients/\allowbreak{}BINDING.json and the six complete compressed integer payloads. They are not parameters that may depend on the unknowns. In particular



\begin{verbatim}
beta=2^1064921-1, cone_A=2^1113, cone_B=2^666104.
\end{verbatim}

MF is already shifted by beta, as in the calibrated source34 family. The complete printed-seed recipe and machine/\allowbreak{}tile data accompany the payloads. Coefficient bits and circuit operations are separate complexity measures.



All following names are polynomial abbreviations, not additional unknowns:



\begin{verbatim}
J=6j+1, q=beta*J+1, W=12w+4,
X=qW, Y=(2v+1)q^3, E0=XY,
a=E0+Y, H=4a+3, Delta=a^2+H=(a+2)^2-1,
F=(K0+W)C-z(q-1),
R=(q(q-F)-C)(q^2-1)+(MC+q*MF)J,
k=R+1+hE0, eta=c-Yk,
Derr=(cone_A*F-q)^2+(cone_B*C-q)^2,
Budget=k^2(q^2-Derr)-4q^2*eta^2,
D=3Budget-4, L=XY^2*k, Dt=1+L(L+k),
Dg=1+D*Dt*s^2, Pg=Dt*Dg,
I=r*c^2, Df=1+Delta*I^2,
Q=Delta^2*I^2, delta=Q-1,
V=Df(c-R+4IT)-c, Cy=QV^2-1, Da=delta*Cy,
K=X(3c-2X)+2,
zc=HT+K, Ac=(zc^2+3H^2)(zc^3-19H^3),
S=srH, Z=S(n+QV)-Ac,
u=9S^2*Df, b=S^2*Da, e=81S^2*Pg.
\end{verbatim}

The equation is the single ordinary integer-polynomial equation



\begin{verbatim}
P = [(Z^2-u-b-e)^2-4(ub+ue+be)]^2-64ubeZ^2 = 0.       (1)
\end{verbatim}

The selected literal schedule folds the factor 81 into the guard as Dg81=\allowbreak{}(243Budget-324)\textasteriskcentered{}\allowbreak{}Dt\textasteriskcentered{}\allowbreak{}s\textasciicircum{}\allowbreak{}2+\allowbreak{}81 and computes e=\allowbreak{}S\textasciicircum{}\allowbreak{}2\textasteriskcentered{}\allowbreak{}Dt\textasteriskcentered{}\allowbreak{}Dg81. This is exactly the displayed e and saves a scalar multiplication. The direct companion changes only Z to Sn-Ac. All powers and scalar multiplications are inlined into counted binary +\allowbreak{},-,\textasteriskcentered{}\allowbreak{} operations in the literal graphs. There are no radicals, divisions, implicit witnesses, variable exponents, or unresolved coefficient ports in those graphs.



The standard universal identity is



\begin{verbatim}
P=product over all eps in {-1,1}^3 of
    (Z-eps1*sqrt(u)-eps2*sqrt(b)-eps3*sqrt(e)).          (2)
\end{verbatim}

This is an identity of ordinary polynomials, including dependent, zero, or negative radicands. It follows by successive quadratic norms in a free multiquadratic algebra. Formula (1), not (2), defines the executable polynomial.



\subsection*{2. Unconditional facts and the exact nonzero bootstrap}



J is odd and nonzero. The authentic beta is odd and is 1 modulo 3. Therefore q is even, q=\allowbreak{}2 modulo 3, and |q|\textgreater{}\allowbreak{}=\allowbreak{}beta-1. W is nonzero and is 1 modulo 3, so |X|\textgreater{}\allowbreak{}=\allowbreak{}4|q| and X=\allowbreak{}2 modulo 3. Y has an odd, nonzero multiplier of q\textasciicircum{}\allowbreak{}3. In particular |X|\textgreater{}\allowbreak{}=\allowbreak{}64, |Y|\textgreater{}\allowbreak{}=\allowbreak{}4096, |a|=\allowbreak{}|(X+\allowbreak{}1)Y| is large, H is odd and nonzero, and Delta\textgreater{}\allowbreak{}1.



Writing ell=\allowbreak{}XY\textasciicircum{}\allowbreak{}2 gives Dt=\allowbreak{}1+\allowbreak{}ell(ell+\allowbreak{}1)k\textasciicircum{}\allowbreak{}2\textgreater{}\allowbreak{}=\allowbreak{}1. Df\textgreater{}\allowbreak{}=\allowbreak{}1. Q is either zero, when I=\allowbreak{}0, or at least Delta\textasciicircum{}\allowbreak{}2\textgreater{}\allowbreak{}1; hence delta never vanishes on Z\textasciicircum{}\allowbreak{}11.



Define pi(x)=\allowbreak{}(x\textasciicircum{}\allowbreak{}2+\allowbreak{}3)(x\textasciicircum{}\allowbreak{}3-19). This polynomial has no rational zero. Since



\begin{verbatim}
Ac=H^5*pi(T+K/H),
\end{verbatim}

we have Ac!=\allowbreak{}0 for every integer assignment. If S=\allowbreak{}0, all three scaled radicands vanish and (1) is exactly Ac\textasciicircum{}\allowbreak{}8, which is nonzero. Every zero therefore has s!=\allowbreak{}0 and r!=\allowbreak{}0. This argument deliberately does not yet assert c!=\allowbreak{}0.



Da is positive before extracting roots. If c=\allowbreak{}0, then I=\allowbreak{}Q=\allowbreak{}0 and delta=\allowbreak{}Cy=\allowbreak{}-1, so Da=\allowbreak{}1. If c!=\allowbreak{}0, then I=\allowbreak{}rc\textasciicircum{}\allowbreak{}2!=\allowbreak{}0, |I|\textgreater{}\allowbreak{}=\allowbreak{}c\textasciicircum{}\allowbreak{}2, Df\textgreater{}\allowbreak{}|c|+\allowbreak{}1, and Q\textgreater{}\allowbreak{}1. For integer b0=\allowbreak{}c-R+\allowbreak{}4IT, either b0=\allowbreak{}0, giving |V|=\allowbreak{}|c|\textgreater{}\allowbreak{}=\allowbreak{}1, or |V|\textgreater{}\allowbreak{}=\allowbreak{}Df-|c|\textgreater{}\allowbreak{}1. Thus Cy\textgreater{}\allowbreak{}0, delta\textgreater{}\allowbreak{}0, and Da\textgreater{}\allowbreak{}0. Together Df\textgreater{}\allowbreak{}0, Dt\textgreater{}\allowbreak{}0, and Da\textgreater{}\allowbreak{}0 hold in both cases.



At a zero a factor in (2) vanishes over C. If Dg\textless{}\allowbreak{}0, the Pg term is the only nonzero imaginary term. This is impossible since S!=\allowbreak{}0. Thus Dg\textgreater{}\allowbreak{}=\allowbreak{}0. If Budget\textless{}\allowbreak{}=\allowbreak{}0, then D\textless{}\allowbreak{}=\allowbreak{}-4 and, using Dt\textgreater{}\allowbreak{}=\allowbreak{}1 and s\textasciicircum{}\allowbreak{}2\textgreater{}\allowbreak{}=\allowbreak{}1,



\begin{verbatim}
Dg=1+D*Dt*s^2<=-3,
\end{verbatim}

a contradiction. Hence Budget\textgreater{}\allowbreak{}0. The fixed cone coefficients and q are even; consequently both terms of Derr and q\textasciicircum{}\allowbreak{}2 are divisible by 4, and Budget is divisible by 4. Thus Budget\textgreater{}\allowbreak{}=\allowbreak{}4, D\textgreater{}\allowbreak{}=\allowbreak{}8, and in fact Dg\textgreater{}\allowbreak{}0. This explicitly excludes a zero guard radicand.



Budget\textgreater{}\allowbreak{}0 gives k!=\allowbreak{}0, Derr\textless{}\allowbreak{}q\textasciicircum{}\allowbreak{}2, and |eta/\allowbreak{}k|\textless{}\allowbreak{}1/\allowbreak{}2. If c=\allowbreak{}0, eta=\allowbreak{}-Yk and



\begin{verbatim}
Budget=k^2[q^2-Derr-4q^2Y^2]<=0,
\end{verbatim}

contradiction. Therefore c!=\allowbreak{}0, I!=\allowbreak{}0, Q\textgreater{}\allowbreak{}1, and all three normalized radicands are now strictly positive.



The disk inequality Derr\textless{}\allowbreak{}q\textasciicircum{}\allowbreak{}2 implies



\begin{verbatim}
0<cone_A*F/q<2, 0<cone_B*C/q<2.
\end{verbatim}

Since both cone coefficients are at least 8, F,C,q-F-C have q's common nonzero sign. Put m=\allowbreak{}|Y| and k0=\allowbreak{}|k|. The centered strip supplies



\begin{verbatim}
k0(m-1/2)<|c|<k0(m+1/2), sign(c)=sign(kY).             (3)
\end{verbatim}

No premature positive claim about eta or k-eta has been made.



\subsection*{3. Prime-three separation and coprime splitting}



Modulo 3, X=\allowbreak{}-1, so a=\allowbreak{}H=\allowbreak{}Delta=\allowbreak{}0. Since ell=\allowbreak{}XY\textasciicircum{}\allowbreak{}2 is either 0 or -1 modulo 3, ell(ell+\allowbreak{}1)=\allowbreak{}0. Hence



\begin{verbatim}
Dt=Df=1, Q=0, delta=Cy=-1, Da=1 (mod 3).              (4)
\end{verbatim}

Divide a vanishing factor by S. It says that a signed sum of



\begin{verbatim}
3sqrt(Df), sqrt(Da), 9sqrt(Pg)
\end{verbatim}

is rational. The squared 3-adic orders of these terms are respectively



\begin{verbatim}
2, 0, 4+v3(Pg).
\end{verbatim}

All radicands are positive integers. In a nontrivial squareclass their square roots are rational multiples of one common irrational root. If two or more terms belonged to that class, their rational coefficients would have distinct 3-adic orders and a unique least order; their signed sum could not vanish. A singleton class cannot cancel either. Linear independence of distinct squarefree radicals therefore forces all three roots to be rational, and thus integral. This argument does not require an even I or an odd guard radicand.



We obtain integers f, ytilde, gprod with



\begin{verbatim}
f^2=Df, ytilde^2=Da, gprod^2=Pg.
\end{verbatim}

The essential congruence Dg=\allowbreak{}1 modulo Dt gives gcd(Dt,Dg)=\allowbreak{}1. Both factors are positive, so unique prime factorization splits their square product:



\begin{verbatim}
Dt=tau^2, Dg=g^2                                     (5)
\end{verbatim}

for positive integers tau,g. This is exactly the selective coprime fusion, not an assertion that a general product square splits.



The rational y=\allowbreak{}ytilde/\allowbreak{}delta satisfies



\begin{verbatim}
QV^2-(Q-1)y^2=1.                                    (6)
\end{verbatim}

Adjacent-coefficient Pell descent makes y integral. For completeness, if y=\allowbreak{}A/\allowbreak{}B in lowest terms then B\textasciicircum{}\allowbreak{}2 divides Q-1. The transformations



\begin{verbatim}
V'=(2Q-1)V-2(Q-1)y,
y'=(2Q-1)y-2QV
\end{verbatim}

preserve the norm and the integrality of V'; their matrix is unimodular over Z, so they preserve the least common denominator of the pair, namely B. After reflecting signs take V,y\textgreater{}\allowbreak{}0. When V\textgreater{}\allowbreak{}1, the norm yields V\textless{}\allowbreak{}y\textless{}\allowbreak{}QV/\allowbreak{}(Q-1), which makes -V\textless{}\allowbreak{}V'\textless{}\allowbreak{}V; the norm excludes V'=\allowbreak{}0. Descent reaches |V|=\allowbreak{}1, where |y|=\allowbreak{}1, so B=\allowbreak{}1.



Thus the factor equation is an integer identity and implies S|Ac. Since H divides S, reducing the carrier modulo H gives exactly



\begin{verbatim}
H | K^5.                                            (7)
\end{verbatim}

We do not incorrectly infer H|K from (7) alone.



\subsection*{4. Signed folding and ranks before projection}



The remaining argument uses precisely the centered signed source retained in source/\allowbreak{}centered11-proof.md, Sections 5.2-5.4, and absolute folding from source/\allowbreak{}signed-quadratic-proof.md, Section 2. We state their hypotheses and explain where they were established, to avoid hiding a positivity or parity assumption.



Put B0=\allowbreak{}|a+\allowbreak{}2|, C0=\allowbreak{}|c|, I0=\allowbreak{}|I|. We have B0\textgreater{}\allowbreak{}=\allowbreak{}3, 0\textless{}\allowbreak{}C0\textless{}\allowbreak{}=\allowbreak{}I0, C0\textasciicircum{}\allowbreak{}2|I0, and



\begin{verbatim}
f^2-(B0^2-1)I0^2=1, |f|>2I0,
rho=Delta*I0, Q=rho^2,
rho^2 V^2-(rho^2-1)y^2=1,
V=-c (mod |f|), V=-R (mod C0).                       (8)
\end{verbatim}

The congruences follow from the displayed definition of V and the exact identity



\begin{verbatim}
V+R=Delta I^2(c-R)+4I Df T.
\end{verbatim}

Classify the strong pair by f\_\allowbreak{}abs=\allowbreak{}chi\_\allowbreak{}B0(m0), I0=\allowbreak{}psi\_\allowbreak{}B0(m0). Classify the auxiliary pair by rho|V|=\allowbreak{}chi\_\allowbreak{}rho(nu), |y|=\allowbreak{}psi\_\allowbreak{}rho(nu). Its index nu is odd, since chi\_\allowbreak{}rho at an even index is plus or minus 1 modulo rho. The odd quotient-polynomial identity and the first congruence in (8) give psi\_\allowbreak{}B0(nu)=\allowbreak{}plus or minus C0 modulo |f|. Fold nu modulo 2m0 and reflect to a p with 0\textless{}\allowbreak{}=\allowbreak{}p\textless{}\allowbreak{}=\allowbreak{}m0. Since psi\_\allowbreak{}B0(p)+\allowbreak{}C0\textless{}\allowbreak{}=\allowbreak{}2I0\textless{}\allowbreak{}|f|, no wrapping is possible, and C0=\allowbreak{}psi\_\allowbreak{}B0(p), p\textgreater{}\allowbreak{}=\allowbreak{}1. Thus the omitted main norm is genuinely recovered before projection.



The unchanged common-sign packing bound is



\begin{verbatim}
|q|(|q|^2-1)<R<|q|^4-2.
\end{verbatim}

The unchanged first norm, the cubic Y scale, k=\allowbreak{}R+\allowbreak{}1+\allowbreak{}hXY, and the centered strip (3) give the pre-projection first-index bounds in the retained centered proof. Its strong divisibility argument needs C0\textasciicircum{}\allowbreak{}2|I0, not an odd multiplier: writing m0=\allowbreak{}pv and D0=\allowbreak{}chi\_\allowbreak{}B0(p), composition gives I0=\allowbreak{}C0\textasteriskcentered{}\allowbreak{}psi\_\allowbreak{}D0(v), while psi\_\allowbreak{}D0(v)=\allowbreak{}vD0\textasciicircum{}\allowbreak{}(v-1) modulo C0 and gcd(D0,C0)=\allowbreak{}1. Therefore C0|v and C0|m0. The auxiliary congruences then give R=\allowbreak{}plus or minus p modulo C0. The strict inherited bounds 0\textless{}\allowbreak{}R,p\textless{}\allowbreak{}C0/\allowbreak{}2 force p=\allowbreak{}R.



These steps precede any use of (7). Their conclusions are



\begin{verbatim}
R=2r0+1, r0>=24,
c=sign(c)*psi_(a+2)(R),
|Y|>(|X|-3)^r0/2-1.                                 (9)
\end{verbatim}

The last ratio estimate follows from the same absolute centered strip, so replacing the old determinant guard by Budget has not weakened its hypotheses. No sign of X,Y,a,c,I,V,H has been discarded. As in the inherited argument, (9) implies, for x=\allowbreak{}|X|,



\begin{verbatim}
|H|>x(x-3)^r0.                                      (10)
\end{verbatim}

Indeed |H|\textgreater{}\allowbreak{}=\allowbreak{}4(x-1)|Y|-3, and substitution from (9) reduces the required strict inequality to (x-2)(x-3)\textasciicircum{}\allowbreak{}r0\textgreater{}\allowbreak{}4x-1, valid here.



\subsection*{5. Powered divisibility restores exact projection}



Put Kquad=\allowbreak{}X(X+\allowbreak{}2ac)-1 and sigma=\allowbreak{}sign(c). The universal identity is



\begin{verbatim}
2Kquad=HcX-K.                                       (11)
\end{verbatim}

Since H is odd, (7) implies H|Kquad\textasciicircum{}\allowbreak{}5. The signed Pell projection identity is



\begin{verbatim}
2^R*Kquad=Pi (mod H),
Pi=(sigma X-2^R)(2^R sigma X+1).                     (12)
\end{verbatim}

It is the exact congruence from source/\allowbreak{}signed-quadratic-proof.md, not a statement assuming H|Kquad. Raising (12) to the fifth power gives H|Pi\textasciicircum{}\allowbreak{}5. There is no cancellation of a nonunit modulo a composite H.



The actual beta implies x\textgreater{}\allowbreak{}=\allowbreak{}4(beta-1)\textgreater{}\allowbreak{}=\allowbreak{}2\textasciicircum{}\allowbreak{}64. Each factor of Pi has absolute value smaller than 2\textasciicircum{}\allowbreak{}(R+\allowbreak{}1)x, so



\begin{verbatim}
|Pi|^5<2^(10R+10)x^10=2^(20r0+20)x^10.
\end{verbatim}

Meanwhile (10) and x-3\textgreater{}\allowbreak{}x/\allowbreak{}2 give |H|\textgreater{}\allowbreak{}x\textasciicircum{}\allowbreak{}(r0+\allowbreak{}1)/\allowbreak{}2\textasciicircum{}\allowbreak{}r0. For r0\textgreater{}\allowbreak{}=\allowbreak{}24 and x\textgreater{}\allowbreak{}=\allowbreak{}2\textasciicircum{}\allowbreak{}64,



\begin{verbatim}
x^(r0-9)>=2^(64(r0-9))>2^(21r0+20),
\end{verbatim}

because 64(r0-9)-(21r0+\allowbreak{}20)=\allowbreak{}43r0-596\textgreater{}\allowbreak{}=\allowbreak{}436. Hence |Pi|\textasciicircum{}\allowbreak{}5\textless{}\allowbreak{}|H|. The integer divisibility forces Pi=\allowbreak{}0. Its second factor is odd and nonzero, so sigma X=\allowbreak{}2\textasciicircum{}\allowbreak{}R.



Now define d=\allowbreak{}sigma\textasteriskcentered{}\allowbreak{}chi\_\allowbreak{}(a+\allowbreak{}2)(R). The signed Pell projection gives H|(d-X-ac), restoring the omitted main root and exact main congruence. The centered orientation and uniform half-binomial argument in source/\allowbreak{}centered11-proof.md Sections 5.2-5.3 now yield X,c,k,Y\textgreater{}\allowbreak{}0 and 0\textless{}\allowbreak{}eta/\allowbreak{}k\textless{}\allowbreak{}1/\allowbreak{}4. Section 5.4 explicitly regenerates an odd strong/\allowbreak{}auxiliary pair and an actual old signed Q41 witness. This is an existential regeneration after rank and interval recovery, not a false odd normalization of the incoming r.



Every zero of (1) therefore yields a finite syntactic ZFC refutation under the unchanged named ordinary source34/\allowbreak{}compiler calibration.



\subsection*{6. Supported-modulus interpolation with a fixed strong pair}



For clarity the needed stronger interpolation theorem is restated, including the argument beyond the original modulus. Its complete proof also appears in source/\allowbreak{}POWERED\_\allowbreak{}SATURATION\_\allowbreak{}PARENT.md, Section 5.



Let A0\textgreater{}\allowbreak{}1 be even, Delta=\allowbreak{}A0\textasciicircum{}\allowbreak{}2-1, R=\allowbreak{}3 modulo 4, and c=\allowbreak{}psi\_\allowbreak{}A0(R)\textgreater{}\allowbreak{}0 with c=\allowbreak{}R modulo 4. For a prescribed positive s, the finite Pell-matrix period supplies an even positive m0 and a strong pair



\begin{verbatim}
I=psi_A0(m0), f=chi_A0(m0), 4s c^2 | I.
\end{verbatim}

Keep this pair fixed throughout. The inherited valuation lemma gives v\_\allowbreak{}p(m0)\textless{}\allowbreak{}=\allowbreak{}v\_\allowbreak{}p(I) for every p|I. The exact squared-modulus period lambda=\allowbreak{}4m0 f is a period modulo f\textasciicircum{}\allowbreak{}2; the generally insufficient 4m0 is not used. Original full-residue interpolation supplies one positive nu0=\allowbreak{}3 modulo 4 and T0 with



\begin{verbatim}
V_nu0=f^2(c-R+4IT0)-c,
V_nu=chi_rho(nu)/rho, rho=Delta I,
\end{verbatim}

at positive odd nu. Set d0=\allowbreak{}lcm(lambda,4I). For p|I,



\begin{verbatim}
v_p(d0)=v_p(4I),                                    (13)
\end{verbatim}

because gcd(f,I)=\allowbreak{}1 and v\_\allowbreak{}p(m0)\textless{}\allowbreak{}=\allowbreak{}v\_\allowbreak{}p(I). For h\textgreater{}\allowbreak{}=\allowbreak{}0 put nu(h)=\allowbreak{}nu0+\allowbreak{}d0\textasteriskcentered{}\allowbreak{}h and



\begin{verbatim}
T(h)=[V_nu(h)+c-f^2(c-R)]/(4I f^2).
\end{verbatim}

These values are integral. The index shift preserves the required residue modulo f\textasciicircum{}\allowbreak{}2 through lambda; it preserves divisibility by 4I through nu=\allowbreak{}nu0 modulo 4I and the identity V\_\allowbreak{}nu=\allowbreak{}-nu modulo 4rho\textasciicircum{}\allowbreak{}2 at nu=\allowbreak{}3 modulo 4. Since gcd(f\textasciicircum{}\allowbreak{}2,4I)=\allowbreak{}1, the product denominator divides the numerator. Each h also has integral auxiliary ordinate psi\_\allowbreak{}rho(nu(h)).



For every p|I and distinct h1,h2,



\begin{verbatim}
v_p(T(h1)-T(h2))=v_p(h1-h2).                        (14)
\end{verbatim}

Here are the lifting details. If D0 is a positive multiple of 4, Pell addition yields



\begin{verbatim}
V_(nu+D0)-V_nu
  =V_nu[chi_rho(D0)-1]
   +psi_rho(nu)(rho^2-1)*psi_rho(D0)/rho.
\end{verbatim}

For p|rho and odd nu, psi\_\allowbreak{}rho(nu) and rho\textasciicircum{}\allowbreak{}2-1 are p-adic units. The elementary lifting formulas are



\begin{verbatim}
v_p(psi_rho(D0)/rho)=v_p(D0),
v_p(chi_rho(D0)-1)=2(v_p(rho)+v_p(D0))   for odd p,
v_2(chi_rho(D0)-1)=2(v_2(rho)+v_2(D0))-1.
\end{verbatim}

They follow from psi\_\allowbreak{}rho(4)=\allowbreak{}4rho(2rho\textasciicircum{}\allowbreak{}2-1), chi\_\allowbreak{}rho(4)-1=\allowbreak{}8rho\textasciicircum{}\allowbreak{}2(rho\textasciicircum{}\allowbreak{}2-1), binomial expansion of the norm-one unit, and the doubling formula. For odd p, multiplying the index by p raises the second-coordinate valuation exactly once; multiplying by a p-adic unit preserves it. The last binomial term at p=\allowbreak{}3 and v\_\allowbreak{}p(rho)=\allowbreak{}1 has valuation 3\textgreater{}\allowbreak{}2, so the smallest exceptional case is also strict. At p=\allowbreak{}2, odd index multiplication preserves v2(psi\_\allowbreak{}rho(4)), and each doubling raises it exactly once because the even-index chi is odd. The first-coordinate formulas then follow from (chi-1)(chi+\allowbreak{}1)=\allowbreak{}(rho\textasciicircum{}\allowbreak{}2-1)psi\textasciicircum{}\allowbreak{}2; chi+\allowbreak{}1 is a p-adic unit for odd p and has 2-adic valuation 1.



Thus the two summands in the displayed addition formula have distinct valuations, and their sum has valuation v\_\allowbreak{}p(D0). Substitute D0=\allowbreak{}d0\textasteriskcentered{}\allowbreak{}|h1-h2|, use (13), and divide by 4I f\textasciicircum{}\allowbreak{}2 to get (14). Reversing the indices covers negative increments.



For every positive modulus M whose prime divisors all divide I, (14) makes h-\textgreater{}\allowbreak{}T(h) a well-defined injective, hence bijective map on each finite ring Z/\allowbreak{}p\textasciicircum{}\allowbreak{}eZ for p\textasciicircum{}\allowbreak{}e|M. Finite CRT chooses one h realizing any prescribed residue of T modulo M; take a nonnegative representative. The same strong pair I,f therefore allows every residue modulo every such M, even when M does not divide I. This stronger result is exactly what permits dropping s from I while retaining sr in S.



\subsection*{7. Completeness and its quantifier order}



Start with a finite syntactic ZFC refutation and the canonical centered outer witness supplied by the inherited calibration. It fixes j,C,z,h,v,w,c and has



\begin{verbatim}
A0=a+2>1 even, R=3 (mod 4), c=psi_A0(R)>0,
c=R (mod 4), Dt=tau^2,
Derr<=q^2/2, 0<eta/k<1/4,
Budget>q^2 k^2/4>=64.
\end{verbatim}

These are the canonical centered-budget estimates, not claims about arbitrary old determinant witnesses. In particular D=\allowbreak{}3Budget-4\textgreater{}\allowbreak{}0. By (4), Dt=\allowbreak{}1 modulo 3, so D\textasteriskcentered{}\allowbreak{}Dt=\allowbreak{}2 modulo 3 and is a positive nonsquare. Ordinary Pell existence supplies positive s,g with



\begin{verbatim}
g^2-D*Dt*s^2=1.
\end{verbatim}

This is exactly our Dg=\allowbreak{}g\textasciicircum{}\allowbreak{}2, without a forced even ordinate. The coprime product is Pg=\allowbreak{}(tau\textasteriskcentered{}\allowbreak{}g)\textasciicircum{}\allowbreak{}2.



Now choose the single strong pair I,f of Section 6 with 4s c\textasciicircum{}\allowbreak{}2|I, and set



\begin{verbatim}
r=I/c^2, M=sr.
\end{verbatim}

Every prime factor of M divides I: those of r do because r|I, and those of s do because s|I. No claim M|I is made. The canonical main projection gives H|Kquad, hence H|K by (11); put kappa=\allowbreak{}K/\allowbreak{}H in Z.



The classical intersective quintic pi(x)=\allowbreak{}(x\textasciicircum{}\allowbreak{}2+\allowbreak{}3)(x\textasciicircum{}\allowbreak{}3-19) has a root modulo every positive integer. Away from 2,3,19, either its cubic has a simple root, or irreducible-cubic Frobenius is a 3-cycle, making its discriminant -27\textasteriskcentered{}\allowbreak{}19\textasciicircum{}\allowbreak{}2 a square and giving a simple root of the quadratic. At 19 the quadratic root 4 is simple. At 2 the cubic root 1 is simple. At 3, x=\allowbreak{}1+\allowbreak{}3u reduces x\textasciicircum{}\allowbreak{}3=\allowbreak{}19 to u+\allowbreak{}3u\textasciicircum{}\allowbreak{}2+\allowbreak{}3u\textasciicircum{}\allowbreak{}3=\allowbreak{}2, whose derivative is 1 modulo 3. Lifting and finite CRT finish the argument.



Choose a residue a0 modulo M such that pi(a0+\allowbreak{}kappa)=\allowbreak{}0 modulo M. Supported-modulus interpolation gives T=\allowbreak{}a0 modulo M and an integer auxiliary y, while keeping I,f,r unchanged. Consequently



\begin{verbatim}
Ac=H^5*pi(T+kappa),
Ac/S=H^4*pi(T+kappa)/(sr) in Z.
\end{verbatim}

Finally choose the unrestricted integer



\begin{verbatim}
n=Ac/S+3f+delta*y+9tau*g-QV.
\end{verbatim}

A factor in (2) vanishes. For the direct companion omit -QV. All eleven coordinates are integers; no sign restriction is imposed on n or T.



The witness order is canonical outer data, guard Pell s/\allowbreak{}g, one fixed strong pair I/\allowbreak{}f/\allowbreak{}r, one intersective residue modulo sr, T/\allowbreak{}y, then n. No strong pair changes after the residue is chosen. Both zero-existence directions are proved, giving the stated consistency interpretation under the unchanged ordinary calibration.



\subsection*{8. Exact degree: cancellation, not a syntactic estimate}



All fixed coefficients have degree zero. The relevant total degrees are



\begin{verbatim}
I:3, S:8, Delta:12, Df:18, Q:30, V:23,
Cy:76, Da:106, Dt:34, Budget:24, Dg:60, Pg:94,
Ac:35, u:34, b:122, e:110.
\end{verbatim}

For the centered equation write B=\allowbreak{}Sn-Ac and Lc=\allowbreak{}SQV, so degB=\allowbreak{}35, degLc=\allowbreak{}61, and Z=\allowbreak{}B+\allowbreak{}Lc. The exact cancellation identity is



\begin{verbatim}
omega=Z^2-b
  =B^2+2LcB+S^2[Q(V^2+1)-1].                       (15)
\end{verbatim}

The three terms have degrees at most 70,96,92. Thus omega has exact degree 96. Rewrite the inner bracket of (1) as



\begin{verbatim}
E8=(omega-u-e)^2-4[ue+b(u+e)].                       (16)
\end{verbatim}

The term -4be has degree 232; all competitors have degree at most 220, including e\textasciicircum{}\allowbreak{}2, and hence it is the unique highest term. The correction 64ubeZ\textasciicircum{}\allowbreak{}2 has degree 34+\allowbreak{}122+\allowbreak{}110+\allowbreak{}122=\allowbreak{}388, below 464. Therefore



\begin{verbatim}
degP=464, P_top=16*b_top^2*e_top^2.                  (17)
\end{verbatim}

For the direct equation degZ=\allowbreak{}35 and b, of degree 122, is the unique largest squared argument. Its norm has exact degree 488 and top form b\_\allowbreak{}top\textasciicircum{}\allowbreak{}4.



Here is a fully attained authentic coefficient for the centered equation. Put



\begin{verbatim}
q0=6beta*j, W0=12w, X0=q0W0, Y0=2v*q0^3,
a0=X0Y0, F0=12wC-6beta*jz, R0=-q0^3F0,
k0=hX0Y0, Dt0=h^2X0^4Y0^6,
Budget0=-4q0^2Y0^2k0^2,
Pg0=-12s^2h^6q0^2X0^10Y0^16,
S0=4sra0, Q0=a0^4r^2c^4,
V0=a0^2r^2c^4q0^3F0.
\end{verbatim}

Substitution in (17) gives the full highest homogeneous form



\begin{verbatim}
P_top=2^452*3^322*beta^260
  *j^260*w^52*v^64*h^12*s^12*r^24*c^32
  *(12wC-6beta*jz)^4.                               (18)
\end{verbatim}

Thus



\begin{verbatim}
[j^260 w^56 v^64 h^12 s^12 r^24 c^32 C^4]P
  =2^460*3^326*beta^260>0.
\end{verbatim}

The exponent sum is 464. The coefficient uses only the authentic positive beta; no relation among the other five fixed ports can cancel it. This proves degree after specialization to the actual tuple, not only with generic symbolic coefficients. The literal independent audit binds (15)-(18) to the exported graph and verifies their exact polynomial identities.



\subsection*{9. Source boundary and reproducibility}



The local source directory preserves the complete centered, signed-folding, strong-index, calibration, and prior arithmetic proofs with their notices. COPRIME\_\allowbreak{}PROJECTION\_\allowbreak{}PARENT.md records the selective-fusion parent, while POWERED\_\allowbreak{}SATURATION\_\allowbreak{}PARENT.md records the earlier powered projection and supported-modulus isometry. Our contribution is their compatible eleven-variable composition, with the nonzero bootstrap and prime-three separation redone after decoupling I and S.



The public comparison snapshot requested for this round is ProveIt commit 45b399ec\allowbreak{}51dfbecb\allowbreak{}891dd51d\allowbreak{}8a9e4b6b\allowbreak{}7c361ac6. The mathematical files inherited through the supplied archives retain their own original provenance pins, principally 84e16e43\allowbreak{}092614ba\allowbreak{}e262ca59\allowbreak{}651cb9d7\allowbreak{}7e31a255 and the coefficient recipe's earlier source commit. A comparison pin does not rewrite an inherited source's provenance. The release ledger records the exact bytes of every retained dependency.



The equation is completely fixed by its eleven variable names, literal graph, six authenticated coefficient payloads, and output node. Reproduction checks include topology and full liveness, independent mathematical-register transcription, universal norm and centering identities, authentic coefficient hashes and recipes, exact degree, signed modular replay, and rejected corruptions. These finite checks certify the arithmetic artifact. The unbounded Pell/\allowbreak{}interpolation/\allowbreak{}source arguments remain ordinary proofs as explicitly stated above.

% Complete proof from norm/DEGREE_PROOF.md; literal text, no Markdown math interpretation.
\section{Exact degrees and literal cost of the round-15 norm pair}



\subsection*{Frozen circuits and domain}



The selected files are \texttt{integral\_\allowbreak{}centered\_\allowbreak{}paired.json} (11 unrestricted integer variables, exact degree 464, 119 live binary operations) and \texttt{integral\_\allowbreak{}direct\_\allowbreak{}paired.json} (11 unrestricted integer variables, exact degree 488, 118 live binary operations). Their mathematical prefix is the frozen three-adic construction in \texttt{../\allowbreak{}scalar/\allowbreak{}PROOF.md}. The complete six ports are beta, K0, MC, MF, cone\_\allowbreak{}A, cone\_\allowbreak{}B; all six authentic byte payloads and their binding are included under \texttt{coefficients/\allowbreak{}}.



Every binary addition, subtraction and multiplication is charged, including scalar multiplication and squaring. Reuse is free. Integer literals, variables, and the six fixed coefficient ports are leaves, not operations and not new unknowns. Every retained arithmetic node is live. The selected pair has respectively (74 multiplications, 24 additions, 21 subtractions) and (74 multiplications, 23 additions, 21 subtractions).



The generator is original local code and imports no upstream executable. The independent skeptic's full register-cut audit is in \texttt{../\allowbreak{}scalar\_\allowbreak{}audit/\allowbreak{}}; \texttt{check\_\allowbreak{}norm.py} independently checks full-byte bindings, all 35 coefficients of the universal radical identity, the center identity, topology, and the exact leading coefficient of every literal DAG node (with the proved cancellation substituted at its one identified node).



\subsection*{A paired quadratic-norm identity}



For formal independent A,B,G,Z put



N(Z;A,B,G) =\allowbreak{} [(Z\textasciicircum{}\allowbreak{}2+\allowbreak{}A-B-G)\textasciicircum{}\allowbreak{}2 - 4(AZ\textasciicircum{}\allowbreak{}2+\allowbreak{}BG)]\textasciicircum{}\allowbreak{}2 - 64(AZ\textasciicircum{}\allowbreak{}2)(BG).



In the free triquadratic algebra this equals



product over eps in \{-1,1\}\textasciicircum{}\allowbreak{}3 of (Z-eps1 sqrt(A)-eps2 sqrt(B)-eps3 sqrt(G)).



Equivalently it is [(Z\textasciicircum{}\allowbreak{}2-A-B-G)\textasciicircum{}\allowbreak{}2-4(AB+\allowbreak{}AG+\allowbreak{}BG)]\textasciicircum{}\allowbreak{}2-64ABGZ\textasciicircum{}\allowbreak{}2. These are identities of complete integer polynomials, not numerical ray checks. Multiplying all eight conjugates gives exactly 35 monomials and zero coefficients on all seven non-scalar radical basis elements. Reusing AZ\textasciicircum{}\allowbreak{}2 and BG saves one binary multiplication over the source's symmetric norm schedule.



In the selected circuits, A=\allowbreak{}9S\textasciicircum{}\allowbreak{}2 Df, B=\allowbreak{}S\textasciicircum{}\allowbreak{}2 delta Cy, G=\allowbreak{}81S\textasciicircum{}\allowbreak{}2 P. The direct coordinate is Z=\allowbreak{}Sn-Ac; the centered coordinate is Z=\allowbreak{}S(n+\allowbreak{}QV)-Ac. The center is an integral shear, so it preserves integer-zero existence without any congruence condition on n.



\subsection*{Weighted guard: one additional exact multiplication saving}



The selected literal graph stores D81=\allowbreak{}243Budget-324, Dg81=\allowbreak{}81+\allowbreak{}D81 Dt s\textasciicircum{}\allowbreak{}2, P81=\allowbreak{}Dt Dg81, and G=\allowbreak{}S\textasciicircum{}\allowbreak{}2 P81. The complete polynomial identity



81+\allowbreak{}(243Budget-324)Dt s\textasciicircum{}\allowbreak{}2 =\allowbreak{} 81[1+\allowbreak{}(3Budget-4)Dt s\textasciicircum{}\allowbreak{}2]



proves Dg81=\allowbreak{}81Dg and P81=\allowbreak{}81P. Thus G remains exactly 81S\textasciicircum{}\allowbreak{}2P, with no zero-set change, new restriction, division, or change in witness ordering. Absorbing the constant 81 into the existing guard constants saves one scalar multiplication. The proof uses conceptual D,Dg,P; their weighted counterparts are the live literal registers. The checker verifies this entire identity, not numerical instances.



\subsection*{Prefix degrees and their leading forms}



A subscript star denotes the highest homogeneous part in the eleven variables, with the six authenticated ports fixed constants. Let



q\textasteriskcentered{}\allowbreak{}=\allowbreak{}6 beta j, X\textasteriskcentered{}\allowbreak{}=\allowbreak{}12w q\textasteriskcentered{}\allowbreak{}, Y\textasteriskcentered{}\allowbreak{}=\allowbreak{}2v (q\textasteriskcentered{}\allowbreak{})\textasciicircum{}\allowbreak{}3, a\textasteriskcentered{}\allowbreak{}=\allowbreak{}X\textasteriskcentered{}\allowbreak{}Y\textasteriskcentered{}\allowbreak{}, F\textasteriskcentered{}\allowbreak{}=\allowbreak{}12w C-6 beta j z, R\textasteriskcentered{}\allowbreak{}=\allowbreak{}-(q\textasteriskcentered{}\allowbreak{})\textasciicircum{}\allowbreak{}3 F\textasteriskcentered{}\allowbreak{}.



Then the exact leading forms needed below are



I\textasteriskcentered{}\allowbreak{}=\allowbreak{}r c\textasciicircum{}\allowbreak{}2, S\textasteriskcentered{}\allowbreak{}=\allowbreak{}4s r a\textasteriskcentered{}\allowbreak{}, Df\textasteriskcentered{}\allowbreak{}=\allowbreak{}(a\textasteriskcentered{}\allowbreak{})\textasciicircum{}\allowbreak{}2 (I\textasteriskcentered{}\allowbreak{})\textasciicircum{}\allowbreak{}2, Q\textasteriskcentered{}\allowbreak{}=\allowbreak{}(a\textasteriskcentered{}\allowbreak{})\textasciicircum{}\allowbreak{}4 (I\textasteriskcentered{}\allowbreak{})\textasciicircum{}\allowbreak{}2, V\textasteriskcentered{}\allowbreak{}=\allowbreak{}-Df\textasteriskcentered{}\allowbreak{} R\textasteriskcentered{}\allowbreak{}=\allowbreak{}(a\textasteriskcentered{}\allowbreak{})\textasciicircum{}\allowbreak{}2 (I\textasteriskcentered{}\allowbreak{})\textasciicircum{}\allowbreak{}2 (q\textasteriskcentered{}\allowbreak{})\textasciicircum{}\allowbreak{}3 F\textasteriskcentered{}\allowbreak{}, Ac\textasteriskcentered{}\allowbreak{}=\allowbreak{}(4a\textasteriskcentered{}\allowbreak{}T)\textasciicircum{}\allowbreak{}5, Dt\textasteriskcentered{}\allowbreak{}=\allowbreak{}h\textasciicircum{}\allowbreak{}2 (a\textasteriskcentered{}\allowbreak{})\textasciicircum{}\allowbreak{}4 (Y\textasteriskcentered{}\allowbreak{})\textasciicircum{}\allowbreak{}2, D\textasteriskcentered{}\allowbreak{}=\allowbreak{}-12(q\textasteriskcentered{}\allowbreak{})\textasciicircum{}\allowbreak{}2 (Y\textasteriskcentered{}\allowbreak{})\textasciicircum{}\allowbreak{}2 h\textasciicircum{}\allowbreak{}2(a\textasteriskcentered{}\allowbreak{})\textasciicircum{}\allowbreak{}2, P\textasteriskcentered{}\allowbreak{}=\allowbreak{}-12s\textasciicircum{}\allowbreak{}2 h\textasciicircum{}\allowbreak{}6(q\textasteriskcentered{}\allowbreak{})\textasciicircum{}\allowbreak{}2(a\textasteriskcentered{}\allowbreak{})\textasciicircum{}\allowbreak{}10(Y\textasteriskcentered{}\allowbreak{})\textasciicircum{}\allowbreak{}6.



In particular the exact degrees are



I 3; S 8; Df 18; Q 30; V 23; Cy 76; delta Cy 106; Dt 34; Dg 60; P 94; Ac 35; A 34; B 122; G 110.



The dominant term of V is -Df R, not the interpolation term 4Df I T. This distinction is essential after decoupling the strong scale.



The two nonzero norm-leading factors are



B\textasteriskcentered{}\allowbreak{}=\allowbreak{}16s\textasciicircum{}\allowbreak{}2 r\textasciicircum{}\allowbreak{}10 c\textasciicircum{}\allowbreak{}16 (a\textasteriskcentered{}\allowbreak{})\textasciicircum{}\allowbreak{}14(q\textasteriskcentered{}\allowbreak{})\textasciicircum{}\allowbreak{}6(F\textasteriskcentered{}\allowbreak{})\textasciicircum{}\allowbreak{}2, G\textasteriskcentered{}\allowbreak{}=\allowbreak{}-2\textasciicircum{}\allowbreak{}6 3\textasciicircum{}\allowbreak{}5 s\textasciicircum{}\allowbreak{}4 r\textasciicircum{}\allowbreak{}2 h\textasciicircum{}\allowbreak{}6(a\textasteriskcentered{}\allowbreak{})\textasciicircum{}\allowbreak{}12(q\textasteriskcentered{}\allowbreak{})\textasciicircum{}\allowbreak{}2(Y\textasteriskcentered{}\allowbreak{})\textasciicircum{}\allowbreak{}6.



Both are nonzero polynomials for the authentic positive beta.



\subsection*{Centered degree 464: global upper bound and exact attainment}



Write W=\allowbreak{}Sn-Ac and C0=\allowbreak{}SQV, so Z=\allowbreak{}W+\allowbreak{}C0. The complete exact identity



Z\textasciicircum{}\allowbreak{}2-B =\allowbreak{} W\textasciicircum{}\allowbreak{}2+\allowbreak{}2SQVW+\allowbreak{}S\textasciicircum{}\allowbreak{}2[Q(V\textasciicircum{}\allowbreak{}2+\allowbreak{}1)-1]



has degree 96: the three displayed terms have degrees 70, 96, and 92. Its degree-96 part is -2S\textasteriskcentered{}\allowbreak{}Q\textasteriskcentered{}\allowbreak{}V\textasteriskcentered{}\allowbreak{}Ac\textasteriskcentered{}\allowbreak{}, which is nonzero. Adding A cannot affect this degree. Hence



v=\allowbreak{}Z\textasciicircum{}\allowbreak{}2+\allowbreak{}A-B-G has degree 110 and top form -G\textasteriskcentered{}\allowbreak{}.



The paired inner bracket is



v\textasciicircum{}\allowbreak{}2-4(AZ\textasciicircum{}\allowbreak{}2+\allowbreak{}BG).



Here v\textasciicircum{}\allowbreak{}2 has degree 220; AZ\textasciicircum{}\allowbreak{}2 has degree 156; BG has degree 232. Therefore the bracket has degree 232 and top form -4B\textasteriskcentered{}\allowbreak{}G\textasteriskcentered{}\allowbreak{}. Its square has degree 464. The final correction 64(AZ\textasciicircum{}\allowbreak{}2)(BG) has degree 388, so it cannot cancel that top form.



Consequently the entire highest homogeneous form is



16(B\textasteriskcentered{}\allowbreak{})\textasciicircum{}\allowbreak{}2(G\textasteriskcentered{}\allowbreak{})\textasciicircum{}\allowbreak{}2 =\allowbreak{}2\textasciicircum{}\allowbreak{}24 3\textasciicircum{}\allowbreak{}10 s\textasciicircum{}\allowbreak{}12 r\textasciicircum{}\allowbreak{}24 c\textasciicircum{}\allowbreak{}32 h\textasciicircum{}\allowbreak{}12(a\textasteriskcentered{}\allowbreak{})\textasciicircum{}\allowbreak{}52(q\textasteriskcentered{}\allowbreak{})\textasciicircum{}\allowbreak{}16(Y\textasteriskcentered{}\allowbreak{})\textasciicircum{}\allowbreak{}12(F\textasteriskcentered{}\allowbreak{})\textasciicircum{}\allowbreak{}4 =\allowbreak{}2\textasciicircum{}\allowbreak{}452 3\textasciicircum{}\allowbreak{}322 beta\textasciicircum{}\allowbreak{}260 j\textasciicircum{}\allowbreak{}260 w\textasciicircum{}\allowbreak{}52 v\textasciicircum{}\allowbreak{}64 h\textasciicircum{}\allowbreak{}12 s\textasciicircum{}\allowbreak{}12 r\textasciicircum{}\allowbreak{}24 c\textasciicircum{}\allowbreak{}32(12wC-6 beta jz)\textasciicircum{}\allowbreak{}4.



For example, the coefficient of



j\textasciicircum{}\allowbreak{}260 w\textasciicircum{}\allowbreak{}56 v\textasciicircum{}\allowbreak{}64 C\textasciicircum{}\allowbreak{}4 h\textasciicircum{}\allowbreak{}12 s\textasciicircum{}\allowbreak{}12 r\textasciicircum{}\allowbreak{}24 c\textasciicircum{}\allowbreak{}32



is exactly 2\textasciicircum{}\allowbreak{}460 3\textasciicircum{}\allowbreak{}326 beta\textasciicircum{}\allowbreak{}260, strictly positive. Its total degree is 464. A second endpoint has j\textasciicircum{}\allowbreak{}264 w\textasciicircum{}\allowbreak{}52 v\textasciicircum{}\allowbreak{}64 z\textasciicircum{}\allowbreak{}4 h\textasciicircum{}\allowbreak{}12 s\textasciicircum{}\allowbreak{}12 r\textasciicircum{}\allowbreak{}24 c\textasciicircum{}\allowbreak{}32 with coefficient 2\textasciicircum{}\allowbreak{}456 3\textasciicircum{}\allowbreak{}326 beta\textasciicircum{}\allowbreak{}264. This proves exact degree over the authentic integers, not just over a finite field and not only along a sampled line.



\subsection*{Direct degree 488}



In the uncentered circuit Z has degree 35, hence Z\textasciicircum{}\allowbreak{}2 has degree 70. The largest term of Z\textasciicircum{}\allowbreak{}2+\allowbreak{}A-B-G is uniquely -B\textasteriskcentered{}\allowbreak{}, degree 122. Its square has degree 244, while AZ\textasciicircum{}\allowbreak{}2 and BG have degrees 104 and 232. Thus the inner bracket has top form (B\textasteriskcentered{}\allowbreak{})\textasciicircum{}\allowbreak{}2, and the output's entire top form is (B\textasteriskcentered{}\allowbreak{})\textasciicircum{}\allowbreak{}4, degree 488. The final correction has degree 336 and cannot cancel it.



One attained monomial is



j\textasciicircum{}\allowbreak{}248 w\textasciicircum{}\allowbreak{}64 v\textasciicircum{}\allowbreak{}56 C\textasciicircum{}\allowbreak{}8 s\textasciicircum{}\allowbreak{}8 r\textasciicircum{}\allowbreak{}40 c\textasciicircum{}\allowbreak{}64,



whose coefficient is exactly 2\textasciicircum{}\allowbreak{}448 3\textasciicircum{}\allowbreak{}312 beta\textasciicircum{}\allowbreak{}248 \textgreater{}\allowbreak{} 0.



\subsection*{Tested but dominated alternatives}



The unfused symmetric comparison graphs cost two additional multiplications: centered 464/\allowbreak{}121 and direct 488/\allowbreak{}120. One saving comes from the paired norm and the other from the weighted guard. Expanding the center cancellation directly gives a syntactic degree-464 circuit but costs 136 operations. Replacing delta\textasteriskcentered{}\allowbreak{}y by y and clearing the auxiliary denominator gives degree 560 and 125 operations uncentered (126 centered). These are not selected. No field-normalization shortcut or unproved division is used in either selected integer circuit.

% Complete proof from scalar/CRT_CARRIER_PROOF.md; literal text, no Markdown math interpretation.
\section{A split scalar carrier lowers the eleven-integer degree to 416}



\subsection*{Status and relationship to the selected primary family}



This is one additional proved degree/\allowbreak{}cost tradeoff, not a replacement of the frozen degree-464 primary. It keeps eleven unrestricted integer variables and the same six actual fixed coefficients. Its new idea is to recover the cubic Y scale through a second, coprime component of the scalar carrier instead of substituting that scale into every arithmetic register. Only the component requiring exact divisibility is raised to the fifth power; the main-Pell projection component retains the previously proved powered-divisibility reduction.



The centered equation has exact degree 416 at 125 live binary operations. Its uncentered companion has exact degree 456 at 124 operations. Their literal schedules are supplied by the separate \texttt{../\allowbreak{}norm\_\allowbreak{}weak/\allowbreak{}} package. The complete ordinary calibration boundary, fixed-tuple payloads, and source proofs are those in PROOF.md and SOURCE\_\allowbreak{}PINS.json. No new formal kernel certificate or minimum claim is made.



\subsection*{1. Complete fixed polynomial}



All eleven variables j,C,z,h,r,v,w,T,c,n,s range over Z. The fixed coefficients beta,K0,MC,MF,cone\_\allowbreak{}A,cone\_\allowbreak{}B are the exact six payloads in coefficients/\allowbreak{}BINDING.json, including shifted MF. Define



\begin{verbatim}
J=6j+1, q=beta*J+1, W=12w+4,
X=qW, Y=qv, E0=XY,
a=E0+Y, H=4a+3, Delta=a^2+H,
F=(K0+W)C-z(q-1),
R=(q(q-F)-C)(q^2-1)+(MC+q*MF)J,
k=R+1+hE0, eta=c-Yk,
Derr=(cone_A*F-q)^2+(cone_B*C-q)^2,
Budget=k^2(q^2-Derr)-4q^2*eta^2,
D=3Budget-4, L=XY^2*k, Dt=1+L(L+k),
Dg=1+D*Dt*s^2, Pg=Dt*Dg,
I=rc^2, Df=1+Delta*I^2,
Q=Delta^2*I^2, delta=Q-1,
V=Df(c-R+4IT)-c, Cy=QV^2-1, Da=delta*Cy,
K=X(3c-2X)+2,
M=2q^2, Dcar=HM,
zc=Dcar*T+M*K+H(v-q^2),
Ac=(zc^2+3Dcar^2)(zc^3-19Dcar^3),
S=srH M^5, Z=S(n+QV)-Ac,
u=9S^2*Df, b=S^2*Da, e=81S^2*Pg.
\end{verbatim}

The final polynomial is exactly the same ordinary three-root norm



\begin{verbatim}
P=[(Z^2-u-b-e)^2-4(ub+ue+be)]^2-64ubeZ^2.             (1)
\end{verbatim}

The direct companion uses Z=\allowbreak{}Sn-Ac. Symbols are polynomial abbreviations, not extra variables. The literal graph may reuse S=\allowbreak{}sr\textasteriskcentered{}\allowbreak{}Dcar\textasteriskcentered{}\allowbreak{}M\textasciicircum{}\allowbreak{}4 and the identity e=\allowbreak{}S\textasciicircum{}\allowbreak{}2\textasteriskcentered{}\allowbreak{}Dt\textasteriskcentered{}\allowbreak{}[(243Budget-324)\textasteriskcentered{}\allowbreak{}Dt\textasteriskcentered{}\allowbreak{}s\textasciicircum{}\allowbreak{}2+\allowbreak{}81]. These are identical polynomials, not new semantic assumptions.



\subsection*{2. The revised bootstrap: restore the scale before excluding c=\allowbreak{}0}



The old argument cannot simply be quoted before recovering the scale: the new Y=\allowbreak{}qv may initially be zero. We explicitly revise the order.



As before q is even, q=\allowbreak{}2 modulo 3, |q|\textgreater{}\allowbreak{}=\allowbreak{}beta-1, W is nonzero, X=\allowbreak{}2 modulo 3, and |X|\textgreater{}\allowbreak{}=\allowbreak{}4|q|. H=\allowbreak{}4a+\allowbreak{}3 is odd and nonzero for every integer a. If v=\allowbreak{}0, a=\allowbreak{}0 and Delta=\allowbreak{}3. If v!=\allowbreak{}0, then



\begin{verbatim}
|a|=|(X+1)qv| >= (|X|-1)|q|,
\end{verbatim}

which is much larger than 4. Hence Delta=\allowbreak{}(a+\allowbreak{}2)\textasciicircum{}\allowbreak{}2-1\textgreater{}\allowbreak{}1 in either case. Dt=\allowbreak{}1+\allowbreak{}ell(ell+\allowbreak{}1)k\textasciicircum{}\allowbreak{}2\textgreater{}\allowbreak{}=\allowbreak{}1 for integer ell=\allowbreak{}XY\textasciicircum{}\allowbreak{}2. Q cannot equal 1, so delta is always nonzero.



M=\allowbreak{}2q\textasciicircum{}\allowbreak{}2 and Dcar=\allowbreak{}HM are nonzero. The same zero-free rational quintic pi(x)=\allowbreak{}(x\textasciicircum{}\allowbreak{}2+\allowbreak{}3)(x\textasciicircum{}\allowbreak{}3-19) gives



\begin{verbatim}
Ac=Dcar^5*pi(zc/Dcar)!=0.
\end{verbatim}

If S=\allowbreak{}0, the three norm radicands vanish and (1) equals Ac\textasciicircum{}\allowbreak{}8. Thus every zero has s!=\allowbreak{}0 and r!=\allowbreak{}0, but c may still be zero at this point.



Df and Da are strictly positive. When c=\allowbreak{}0, I=\allowbreak{}Q=\allowbreak{}0 and Da=\allowbreak{}1. When c!=\allowbreak{}0, r!=\allowbreak{}0 implies I!=\allowbreak{}0, |I|\textgreater{}\allowbreak{}=\allowbreak{}c\textasciicircum{}\allowbreak{}2, and Df\textgreater{}\allowbreak{}|c|+\allowbreak{}1; the same integer-distance argument using V=\allowbreak{}Df(c-R+\allowbreak{}4IT)-c gives V!=\allowbreak{}0, Cy\textgreater{}\allowbreak{}0, and Da\textgreater{}\allowbreak{}0. Therefore if Dg\textless{}\allowbreak{}0 it is the only imaginary radical in a vanishing conjugate factor, a contradiction. Dg\textgreater{}\allowbreak{}=\allowbreak{}0 follows. As before Budget\textless{}\allowbreak{}=\allowbreak{}0 would give Dg\textless{}\allowbreak{}=\allowbreak{}-3; thus Budget\textgreater{}\allowbreak{}0. The parity of q and both cone coefficients gives Budget=\allowbreak{}0 modulo 4, so Budget\textgreater{}\allowbreak{}=\allowbreak{}4 and Dg\textgreater{}\allowbreak{}0.



Modulo 3, X=\allowbreak{}-1 still gives a=\allowbreak{}H=\allowbreak{}Delta=\allowbreak{}0, while ell=\allowbreak{}XY\textasciicircum{}\allowbreak{}2 is 0 or -1. Therefore Dt=\allowbreak{}Df=\allowbreak{}Da=\allowbreak{}1 modulo 3, regardless of v. Prime-three squareclass separation applies to the positive normalized radicals



\begin{verbatim}
3sqrt(Df), sqrt(Da), 9sqrt(Pg).
\end{verbatim}

Their squared orders are 2,0,4+\allowbreak{}v3(Pg). Consequently they are all integers. Coprimality gcd(Dt,Dg)=\allowbreak{}1 splits Pg into Dt=\allowbreak{}tau\textasciicircum{}\allowbreak{}2 and Dg=\allowbreak{}g\textasciicircum{}\allowbreak{}2.



Crucially, the integer factor identity already gives



\begin{verbatim}
S | Ac                                                  (2)
\end{verbatim}

without adjacent-Pell descent and without c!=\allowbreak{}0. The integral norm uses sqrt(Da) directly; no quotient y=\allowbreak{}ytilde/\allowbreak{}delta is needed for this divisibility. This is the point that licenses the new order of proof.



\subsection*{3. The two coprime projections in the one carrier}



We have H=\allowbreak{}3 modulo q and q=\allowbreak{}2 modulo 3, so gcd(H,q)=\allowbreak{}1. H is odd, hence



\begin{verbatim}
gcd(H,M)=1.                                            (3)
\end{verbatim}

Since M\textasciicircum{}\allowbreak{}5 divides S, (2) gives M\textasciicircum{}\allowbreak{}5|Ac. For a prime p with p\textasciicircum{}\allowbreak{}e exactly dividing M, suppose b0=\allowbreak{}v\_\allowbreak{}p(zc)\textless{}\allowbreak{}e. Expand



\begin{verbatim}
Ac=zc^5+3Dcar^2 zc^3-19Dcar^3 zc^2-57Dcar^5.
\end{verbatim}

Dcar is divisible by M. The first term then has valuation 5b0 and each other term has strictly larger valuation, including when p divides 3,19, or 57. That is incompatible with divisibility by p\textasciicircum{}\allowbreak{}(5e). Thus M|zc. Reducing its definition modulo M and using (3) gives



\begin{verbatim}
M | v-q^2.
\end{verbatim}

Put v\_\allowbreak{}old=\allowbreak{}(v-q\textasciicircum{}\allowbreak{}2)/\allowbreak{}(2q\textasciicircum{}\allowbreak{}2) in Z. Then



\begin{verbatim}
v=q^2(2v_old+1), Y=(2v_old+1)q^3.                       (4)
\end{verbatim}

The original nonzero cubic odd-multiplier scale is now restored exactly.



The other component of (2) is H|Ac. Reducing modulo H yields



\begin{verbatim}
H | (MK)^5.
\end{verbatim}

By (3), this implies H|K\textasciicircum{}\allowbreak{}5. Only this weaker main-projection congruence is claimed. The odd-H identity 2Kquad=\allowbreak{}HcX-K and the huge-base argument in PROOF.md Section 5 then recover the exact main projection once the pre-projection ranks are available.



\subsection*{4. Completing soundness in the correct order}



Budget\textgreater{}\allowbreak{}0 supplies k!=\allowbreak{}0 and |c/\allowbreak{}k-Y|\textless{}\allowbreak{}1/\allowbreak{}2. Once (4) is known, |Y|\textgreater{}\allowbreak{}=\allowbreak{}|q|\textasciicircum{}\allowbreak{}3 and Y!=\allowbreak{}0. Therefore c=\allowbreak{}0 would give |Y|\textless{}\allowbreak{}1/\allowbreak{}2, impossible. Now c!=\allowbreak{}0, I!=\allowbreak{}0, Q\textgreater{}\allowbreak{}1, and adjacent-Pell descent is licensed: for the extracted integer ytilde with ytilde\textasciicircum{}\allowbreak{}2=\allowbreak{}Da, y=\allowbreak{}ytilde/\allowbreak{}delta is rational and satisfies



\begin{verbatim}
QV^2-(Q-1)y^2=1.
\end{verbatim}

The descent in PROOF.md Section 3 gives y in Z. All hypotheses of the original signed folding and centered pre-projection rank proof are now present, with v\_\allowbreak{}old in place of the original cubic-scale coordinate. In particular c\textasciicircum{}\allowbreak{}2|I, the two auxiliary residues, the unchanged first norm, the common-sign packing bounds, and the centered strip hold exactly.



Consequently R=\allowbreak{}2r0+\allowbreak{}1 with r0\textgreater{}\allowbreak{}=\allowbreak{}24, the same signed main-Pell classification holds, and |H|\textgreater{}\allowbreak{}|X|(|X|-3)\textasciicircum{}\allowbreak{}r0. Together with H|K\textasciicircum{}\allowbreak{}5 and the authentic |X|\textgreater{}\allowbreak{}=\allowbreak{}2\textasciicircum{}\allowbreak{}64, the powered projection has the same strict exponent margin 436 and restores sign(c)X=\allowbreak{}2\textasciicircum{}\allowbreak{}R. The late centered orientation and half-binomial argument then regenerate an actual old signed Q41 witness and hence a finite syntactic ZFC refutation under the retained ordinary calibration.



The order is S!=\allowbreak{}0, radical realness and Budget positivity, integer-root extraction, carrier divisibility, exact cubic scale, c!=\allowbreak{}0 and adjacent descent, signed ranks, powered main projection, late orientation, source calibration. In particular no large-Y, c!=\allowbreak{}0, or Q\textgreater{}\allowbreak{}1 claim is used before its proof.



\subsection*{5. Completeness and fixed-witness interpolation}



Start with a canonical native outer witness and its original coordinate v\_\allowbreak{}old. Define the new variable



\begin{verbatim}
v=q^2(2v_old+1).
\end{verbatim}

All numerical values of the old front-end registers Y,a,H,Delta,F,R,k,Budget,Dt are then unchanged. Choose the positivity Pell pair s,g for D\textasteriskcentered{}\allowbreak{}Dt as in PROOF.md Section 7. Choose one fixed strong pair I,f with 4s c\textasciicircum{}\allowbreak{}2|I, and set r=\allowbreak{}I/\allowbreak{}c\textasciicircum{}\allowbreak{}2. Supported-modulus interpolation, proved in PROOF.md Section 6, allows every T residue modulo sr with this same strong pair.



The canonical main projection gives H|K; put kappa=\allowbreak{}K/\allowbreak{}H. Also (v-q\textasciicircum{}\allowbreak{}2)/\allowbreak{}M=\allowbreak{}v\_\allowbreak{}old. The new carrier has the exact factorization



\begin{verbatim}
zc=HM(T+kappa+v_old),
Ac=(HM)^5*pi(T+kappa+v_old).
\end{verbatim}

Choose an intersective root for pi(T+\allowbreak{}kappa+\allowbreak{}v\_\allowbreak{}old) modulo sr, then interpolate T in that residue, keeping I,f,r fixed. It follows that



\begin{verbatim}
Ac/S=H^4*pi(T+kappa+v_old)/(sr) in Z.
\end{verbatim}

Finally choose



\begin{verbatim}
n=Ac/S+3f+delta*y+9tau*g-QV
\end{verbatim}

for the centered equation, omitting -QV for the direct one. A conjugate factor vanishes. The witness order is the same as the primary proof; the new coordinate v is fixed with the native outer data, before any Pell or residue choices.



Thus the new equations have an integer zero exactly when the same native source has a finite syntactic refutation. Their universal nonvanishing is ZFC-provably equivalent to Con(ZFC), under exactly the same retained ordinary calibration as the primary equation.



\subsection*{6. Exact degrees and attained coefficients}



The total degrees of the changed prefix are



\begin{verbatim}
Y:2, a:4, Delta:8, k:5, Dt:22, Budget:16, Dg:40, Pg:62,
I:3, Df:14, Q:22, V:19, Cy:60, Da:82,
M:2, Dcar:6, S:16, zc:7, Ac:35,
u:46, b:114, e:94.
\end{verbatim}

For the centered equation let B0=\allowbreak{}Sn-Ac, Lc=\allowbreak{}SQV. Their degrees are 35 and 57. The exact identity



\begin{verbatim}
omega=Z^2-b=B0^2+2LcB0+S^2[Q(V^2+1)-1]
\end{verbatim}

has degree at most 92. Two displayed terms have the same degree, so no unsupported unique-top assertion is made about omega. The paired/\allowbreak{}symmetric norm identity gives the inner bracket



\begin{verbatim}
(omega-u-e)^2-4[ue+b(u+e)].
\end{verbatim}

The term -4be has degree 208; every other term has degree at most 188. The final correction 64ubeZ\textasciicircum{}\allowbreak{}2 has degree 46+\allowbreak{}114+\allowbreak{}94+\allowbreak{}114=\allowbreak{}368, less than 416. Thus the entire highest form is 16b\_\allowbreak{}top\textasciicircum{}\allowbreak{}2 e\_\allowbreak{}top\textasciicircum{}\allowbreak{}2, and its nonzero coefficient below proves exact degree 416.



Write



\begin{verbatim}
q0=6beta*j, X0=12wq0, Y0=q0v, a0=X0Y0,
F0=12wC-q0z, R0=-q0^3F0, k0=R0+h a0.
\end{verbatim}

The tie in k's leading form is essential: R0 is not to be discarded. Directly,



\begin{verbatim}
S0=128sr a0 q0^10,
b_top=2^14*s^2*r^10*c^16*a0^14*q0^26*F0^2,
e_top=-2^16*3^5*s^4*r^2*a0^6*q0^22*Y0^6*k0^6.
\end{verbatim}

Therefore



\begin{verbatim}
P_top=2^64*3^10*s^12*r^24*c^32
               *a0^40*q0^96*Y0^12*k0^12*F0^4.
\end{verbatim}

Select the h\textasciicircum{}\allowbreak{}12 term in k0\textasciicircum{}\allowbreak{}12 and the C\textasciicircum{}\allowbreak{}4 term in F0\textasciicircum{}\allowbreak{}4. This gives



\begin{verbatim}
[j^212 w^56 v^64 h^12 s^12 r^24 c^32 C^4]P
   =2^388*3^278*beta^212 > 0.
\end{verbatim}

The exponent sum is 416. No other term of k0\textasciicircum{}\allowbreak{}12 contributes h\textasciicircum{}\allowbreak{}12, so the tie cannot cancel this coefficient. Its positivity uses only the authentic beta. Selecting instead the R0\textasciicircum{}\allowbreak{}12 term and then the C\textasciicircum{}\allowbreak{}16 term gives the separately actual-DAG-verified endpoint



\begin{verbatim}
[j^224 w^56 v^52 C^16 s^12 r^24 c^32]P
   =2^400*3^290*beta^224 > 0.
\end{verbatim}

Its exponent sum is also 416.



For the direct equation Z has degree 35, and b of degree 114 is the unique largest squared argument. The norm has top form b\_\allowbreak{}top\textasciicircum{}\allowbreak{}4, of degree 456. An attained coefficient is



\begin{verbatim}
[j^216 w^64 v^56 s^8 r^40 c^64 C^8]P_direct
   =2^400*3^280*beta^216 > 0.
\end{verbatim}

The literal norm package independently verifies the cancellation identity and the C\textasciicircum{}\allowbreak{}16 centered and C\textasciicircum{}\allowbreak{}8 direct coefficient certificates against the exported division-free schedules. The finite tests do not replace the signed soundness or completeness argument above.

% Complete proof from norm_weak/DEGREE_PROOF.md; literal text, no Markdown math interpretation.
\section{Weak-Y split-carrier norm: exact degree and cost}



This is a separately audited degree/\allowbreak{}cost tradeoff, not a replacement of the frozen cheaper norm pair. The two graphs use eleven unrestricted integer variables and the same six fully bound authentic ports. The centered graph has exact degree 416 and 125 live binary operations (79 multiplications,24 additions,22 subtractions); its direct sibling has degree 456 and 124 operations (79,23,22).



\subsection*{Definitions and identity boundary}



Retain the complete ordinary prefix and weighted guard of the frozen construction, but use Y=\allowbreak{}qv. Put M=\allowbreak{}2q\textasciicircum{}\allowbreak{}2, Dc=\allowbreak{}HM, M2=\allowbreak{}M\textasciicircum{}\allowbreak{}2, M4=\allowbreak{}M2\textasciicircum{}\allowbreak{}2 and



S=\allowbreak{}sr Dc M4=\allowbreak{}srH M\textasciicircum{}\allowbreak{}5, zc=\allowbreak{}Dc T+\allowbreak{}MK+\allowbreak{}H(v-q\textasciicircum{}\allowbreak{}2), Ac=\allowbreak{}(zc\textasciicircum{}\allowbreak{}2+\allowbreak{}3Dc\textasciicircum{}\allowbreak{}2)(zc\textasciicircum{}\allowbreak{}3-19Dc\textasciicircum{}\allowbreak{}3).



All quantities are polynomial abbreviations, not extra variables. The polynomial remains the full eight-conjugate norm



[(Z\textasciicircum{}\allowbreak{}2+\allowbreak{}A-B-G)\textasciicircum{}\allowbreak{}2-4(AZ\textasciicircum{}\allowbreak{}2+\allowbreak{}BG)]\textasciicircum{}\allowbreak{}2-64(AZ\textasciicircum{}\allowbreak{}2)(BG),



where A=\allowbreak{}9S\textasciicircum{}\allowbreak{}2Df, B=\allowbreak{}S\textasciicircum{}\allowbreak{}2delta Cy, G=\allowbreak{}81S\textasciicircum{}\allowbreak{}2P, and Z=\allowbreak{}S(n+\allowbreak{}QV)-Ac (centered) or Sn-Ac (direct). Its literal guard schedule is D81=\allowbreak{}243Budget-324, Dg81=\allowbreak{}81+\allowbreak{}D81Dt s\textasciicircum{}\allowbreak{}2, P81=\allowbreak{}DtDg81, G=\allowbreak{}S\textasciicircum{}\allowbreak{}2P81. The full polynomial identity Dg81=\allowbreak{}81Dg is checked.



The separate semantic proof must recover M|(v-q\textasciicircum{}\allowbreak{}2) from the carrier BEFORE assuming the old cubic-Y size bounds, c!=\allowbreak{}0 or adjacent-Pell descent. This norm note does not silently inherit those assertions before that step. See the accompanying weak-Y semantic proof for the exact quantifier order.



\subsection*{Highest homogeneous parts}



All degrees count only the eleven variables; every authentic coefficient is fixed. Write



q0=\allowbreak{}6beta j, X0=\allowbreak{}12wq0, Y0=\allowbreak{}q0v, a0=\allowbreak{}X0Y0, F0=\allowbreak{}12wC-q0z, R0=\allowbreak{}-q0\textasciicircum{}\allowbreak{}3F0, k0=\allowbreak{}R0+\allowbreak{}h a0, I0=\allowbreak{}rc\textasciicircum{}\allowbreak{}2, S0=\allowbreak{}128sr a0 q0\textasciicircum{}\allowbreak{}10.



The two terms of k0 both have degree 5; neither may be discarded in the complete top homogeneous form. The exact degrees are



I 3, S 16, Df 14, Q 22, V 19, Cy 60, delta Cy 82, Dt 22, Dg 40, P 62, Ac 35, A 46, B 114, G 94.



The leading forms relevant to the norm are



V0=\allowbreak{}-a0\textasciicircum{}\allowbreak{}2 I0\textasciicircum{}\allowbreak{}2 R0, B0=\allowbreak{}2\textasciicircum{}\allowbreak{}14 s\textasciicircum{}\allowbreak{}2r\textasciicircum{}\allowbreak{}10c\textasciicircum{}\allowbreak{}16 a0\textasciicircum{}\allowbreak{}14 q0\textasciicircum{}\allowbreak{}26 F0\textasciicircum{}\allowbreak{}2, G0=\allowbreak{}-2\textasciicircum{}\allowbreak{}16 3\textasciicircum{}\allowbreak{}5 s\textasciicircum{}\allowbreak{}4r\textasciicircum{}\allowbreak{}2 a0\textasciicircum{}\allowbreak{}6 q0\textasciicircum{}\allowbreak{}22 Y0\textasciicircum{}\allowbreak{}6 k0\textasciicircum{}\allowbreak{}6, Ac0=\allowbreak{}(8a0 q0\textasciicircum{}\allowbreak{}2T)\textasciicircum{}\allowbreak{}5.



All are nonzero over the authentic coefficient tuple.



\subsection*{Exact degree 416 after centering}



Set W=\allowbreak{}Sn-Ac. The universal polynomial identity is



Z\textasciicircum{}\allowbreak{}2-B=\allowbreak{}W\textasciicircum{}\allowbreak{}2+\allowbreak{}2SQVW+\allowbreak{}S\textasciicircum{}\allowbreak{}2[Q(V\textasciicircum{}\allowbreak{}2+\allowbreak{}1)-1].



The displayed degrees are 70,92,92. Thus the expression has degree at most 92. In fact its degree-92 part is -2S0Q0V0Ac0+\allowbreak{}S0\textasciicircum{}\allowbreak{}2Q0V0\textasciicircum{}\allowbreak{}2. The two summands have different C degrees,1 and 2, so this part is nonzero.



Since G has degree 94 and A degree 46, the paired norm's v=\allowbreak{}Z\textasciicircum{}\allowbreak{}2+\allowbreak{}A-B-G has unique top form -G0. In the inner bracket v\textasciicircum{}\allowbreak{}2-4(AZ\textasciicircum{}\allowbreak{}2+\allowbreak{}BG), the degrees are 188,160 and 208, respectively. Hence its top form is -4B0G0. The final correction 64(AZ\textasciicircum{}\allowbreak{}2)(BG) has degree 368, strictly below 416. The complete highest homogeneous form of the final polynomial therefore is



16B0\textasciicircum{}\allowbreak{}2G0\textasciicircum{}\allowbreak{}2 =\allowbreak{}2\textasciicircum{}\allowbreak{}64 3\textasciicircum{}\allowbreak{}10 s\textasciicircum{}\allowbreak{}12r\textasciicircum{}\allowbreak{}24c\textasciicircum{}\allowbreak{}32 a0\textasciicircum{}\allowbreak{}40 q0\textasciicircum{}\allowbreak{}96 Y0\textasciicircum{}\allowbreak{}12 k0\textasciicircum{}\allowbreak{}12 F0\textasciicircum{}\allowbreak{}4.



Taking the maximal-C term in k0\textasciicircum{}\allowbreak{}12 and F0\textasciicircum{}\allowbreak{}4 gives an attained monomial



j\textasciicircum{}\allowbreak{}224 w\textasciicircum{}\allowbreak{}56 v\textasciicircum{}\allowbreak{}52 C\textasciicircum{}\allowbreak{}16 s\textasciicircum{}\allowbreak{}12 r\textasciicircum{}\allowbreak{}24 c\textasciicircum{}\allowbreak{}32



with exact coefficient 2\textasciicircum{}\allowbreak{}400 3\textasciicircum{}\allowbreak{}290 beta\textasciicircum{}\allowbreak{}224\textgreater{}\allowbreak{}0. Its total degree is 416. Another endpoint, choosing the h\textasciicircum{}\allowbreak{}12 branch, is



j\textasciicircum{}\allowbreak{}212 w\textasciicircum{}\allowbreak{}56 v\textasciicircum{}\allowbreak{}64 h\textasciicircum{}\allowbreak{}12 s\textasciicircum{}\allowbreak{}12 r\textasciicircum{}\allowbreak{}24 c\textasciicircum{}\allowbreak{}32 C\textasciicircum{}\allowbreak{}4



with coefficient 2\textasciicircum{}\allowbreak{}388 3\textasciicircum{}\allowbreak{}278 beta\textasciicircum{}\allowbreak{}212\textgreater{}\allowbreak{}0.



Thus the claimed degree is a universal upper bound with an explicit nonzero authentic integer coefficient, not a sampled-ray claim.



\subsection*{Direct sibling: degree 456}



Without the integral shear, deg Z=\allowbreak{}35. In Z\textasciicircum{}\allowbreak{}2+\allowbreak{}A-B-G the unique top form is -B0 (degree 114). Its square dominates BG (degree 208) and AZ\textasciicircum{}\allowbreak{}2 (degree 116), so the final output has whole top form B0\textasciicircum{}\allowbreak{}4 and degree 456. The final correction has degree 324.



The monomial j\textasciicircum{}\allowbreak{}216 w\textasciicircum{}\allowbreak{}64 v\textasciicircum{}\allowbreak{}56 C\textasciicircum{}\allowbreak{}8 s\textasciicircum{}\allowbreak{}8 r\textasciicircum{}\allowbreak{}40 c\textasciicircum{}\allowbreak{}64 has coefficient 2\textasciicircum{}\allowbreak{}400 3\textasciicircum{}\allowbreak{}280 beta\textasciicircum{}\allowbreak{}216\textgreater{}\allowbreak{}0.



\subsection*{Checks and arithmetic cost}



\texttt{check\_\allowbreak{}norm.py} checks all 35 coefficients of the free triquadratic product, the paired/\allowbreak{}symmetric identity, the center identity, the weighted-guard identity, full six-payload byte bindings, graph topology, and exact leading degrees at every actual DAG node with the proved center cancellation substituted at its identified node. Both normal and python-O runs pass. The independent skeptic checks the complete prefix against a separate transcription and exercises corrupted-graph controls.



The generator uses only owned local arithmetic-DAG code. Every binary scalar multiplication and square is charged; reuse and leaves are free. The two candidate files and full payloads are identified by the accompanying release manifest. Frozen lower-cost graphs in \texttt{../\allowbreak{}norm/\allowbreak{}} are unchanged.

% Complete proof from architecture/UNIFIED_PRODUCT_PROOF.md; literal text, no Markdown math interpretation.
\section{One binomial simultaneously certifies a factorial and every CRT modulus}



\subsection*{Scope and novelty}



This is an ordinary, existence-preserving arithmetic replacement of the round14 direct C222 certificate. It uses neither a universal-polynomial black box nor a numerical execution of K215. Its final polynomial is the sum of squares of the entire quadratic residual list in each natural/\allowbreak{}equation.json. Every variable ranges over the nonnegative integers, including zero. All powers and binomials in this proof are explanatory notation: the generator emits the inherited finite ordinary Pell and digit graphs, with no variable-exponent primitive remaining in the polynomial.



The square-output factorial identity and the uniform lower-binomial coefficient bound already occur in the supplied Synchronized Exponents package, which credits the earlier Bootstrap package. They are explicitly inherited, not claimed new. The new architecture identifies the square-radius factorial binomial with the reciprocal CRT product from round14. Their raw ceiling equations then give a new bootstrap inequality V\textless{}\allowbreak{}B\textasciicircum{}\allowbreak{}2, which permits the common radix to depend only on the amplified CRT modulus. There is now one exact binomial, rather than two, and eighteen power requests, rather than nineteen. No additional guard or unproved fixed-point condition is introduced.



The logical endpoint retains the ordinary C222 calculus calibration, the pairwise-coprime Horn-flow theorem, the simultaneous positive-power/\allowbreak{}Pell graph, and the exact natural-domain contraction from the pinned round14 archive. With these retained dependencies, the equation has a natural zero if and only if that calculus has a finite refutation; ZFC proves that its unsolvability is equivalent to Con(ZFC). This is not a proof of consistency or a newly built Lean/\allowbreak{}Rocq/\allowbreak{}object-language ZFC derivation.



\subsection*{1. The unified arithmetic graph}



Let t=\allowbreak{}t0+\allowbreak{}16. F is initially only a proposed natural factorial output. Define



\begin{verbatim}
rho=F^2,        V=(rho+1)^t,
B=binom(rho+t,t),        R=2+t+x+y+F,
K=2R,          N=B^K.
\end{verbatim}

The graph for B is initially an uninterpreted exact-binomial digit graph. Its intended meaning must not be used in the bootstrap below. Impose the same natural ceiling rows as before:



\begin{verbatim}
V+B=BF+u+1,       B=u+v+1.                    (1)
\end{verbatim}

The base radix variant is q=\allowbreak{}N\textasciicircum{}\allowbreak{}K. The squared-base companion is



\begin{verbatim}
q=(N^2)^K.                                    (2)
\end{verbatim}

The square N\textasciicircum{}\allowbreak{}2 already occurs in the representative bound, so the companion adds neither a new variable-power request nor a new exponent index. The companion is useful for subsequent signed-domain compilation because its base is a square. Every natural-domain argument below applies to both variants.



All independent CRT representatives C,A\_\allowbreak{}j are positive. There are nine independent representatives for the derived-label variant and ten for the local-slack variant, with the aggregate equation



\begin{verbatim}
C+sum_j A_j+g=cN^2,       c=9 or c=10.         (3)
\end{verbatim}

The polynomial derived label expression J(A) is NOT included in (3).



\subsection*{2. Uniform coefficient lemma (inherited; complete proof)}



For positive integers r,n and 0\textless{}\allowbreak{}=\allowbreak{}j\textless{}\allowbreak{}=\allowbreak{}n,



\begin{verbatim}
binom(r+n,j)<=(r+1)^n.                       (4)
\end{verbatim}

Induct on n. The case n=\allowbreak{}1 is immediate. The case j=\allowbreak{}0 is immediate for every n. If 1\textless{}\allowbreak{}=\allowbreak{}j\textless{}\allowbreak{}n, Pascal's identity and the induction hypothesis bound the coefficient by 2(r+\allowbreak{}1)\textasciicircum{}\allowbreak{}(n-1)\textless{}\allowbreak{}=\allowbreak{}(r+\allowbreak{}1)\textasciicircum{}\allowbreak{}n. If j=\allowbreak{}n, use



\begin{verbatim}
binom(r+n,n)=product_(i=1)^n (r+i)/i
            <= product_(i=1)^n (r+1).
\end{verbatim}

Unlike a unimodality proof, this requires no hypothesis r\textgreater{}\allowbreak{}=\allowbreak{}n. This is essential because F has not yet been identified as a factorial.



\subsection*{3. Noncircular bootstrap: B\textgreater{}\allowbreak{}F and V\textless{}\allowbreak{}B\textasciicircum{}\allowbreak{}2}



The inherited simultaneous positive-power theorem first recovers V=\allowbreak{}(F\textasciicircum{}\allowbreak{}2+\allowbreak{}1)\textasciicircum{}\allowbreak{}t from an independently positive base and exponent. Thus V\textgreater{}\allowbreak{}=\allowbreak{}1. The second raw row in (1) gives B\textgreater{}\allowbreak{}=\allowbreak{}1. Subtracting the rows gives



\begin{verbatim}
V+v=BF.                                       (5)
\end{verbatim}

Since v is natural, (5) excludes F=\allowbreak{}0 and gives V\textless{}\allowbreak{}=\allowbreak{}BF. Because t\textgreater{}\allowbreak{}=\allowbreak{}1,



\begin{verbatim}
V=(F^2+1)^t>F^2.
\end{verbatim}

Consequently BF\textgreater{}\allowbreak{}F\textasciicircum{}\allowbreak{}2; division by the already positive F gives B\textgreater{}\allowbreak{}F\textgreater{}\allowbreak{}=\allowbreak{}1, so B\textgreater{}\allowbreak{}=\allowbreak{}2, and



\begin{verbatim}
V<=BF<B^2.                                    (6)
\end{verbatim}

These are consequences of raw polynomial rows and an independently valid power. They do not assume that the binomial graph is exact. Now the base B and exponent K are positive, so the inherited power graph recovers N=\allowbreak{}B\textasciicircum{}\allowbreak{}K. The bases N and N\textasciicircum{}\allowbreak{}2 are positive, so the graph recovers the selected q in (2). Since K\textgreater{}\allowbreak{}=\allowbreak{}2,



\begin{verbatim}
q>=N^K=B^(K^2)>=B^4>V.                       (7)
\end{verbatim}

Apply (4) with r=\allowbreak{}F\textasciicircum{}\allowbreak{}2,n=\allowbreak{}t. Every coefficient needed in the sole exact-binomial digit is \textless{}\allowbreak{}=\allowbreak{}V\textless{}\allowbreak{}q. The digit graph therefore really returns



\begin{verbatim}
B=binom(F^2+t,t).
\end{verbatim}

There is no circular factorial, binomial, positivity, or radix assumption in this ordering.



\subsection*{4. Square-output fixed point at the padded horizon}



For completeness the inherited fixed-point theorem is reproved in the restricted t\textgreater{}\allowbreak{}=\allowbreak{}16 range used here. Put H=\allowbreak{}t! and f=\allowbreak{}V/\allowbreak{}B. Equation (1) is equivalent to



\begin{verbatim}
F-1<f<=F,       or F=ceil(f).                  (8)
\end{verbatim}

Indeed 0\textless{}\allowbreak{}=\allowbreak{}u\textless{}\allowbreak{}B, and solving the first row for V gives B(F-1)\textless{}\allowbreak{}V\textless{}\allowbreak{}=\allowbreak{}BF. Conversely, if (8) holds, u=\allowbreak{}V+\allowbreak{}B-BF-1 and v=\allowbreak{}B-u-1 are natural.



The exact product expression is



\begin{verbatim}
f=H product_(i=1)^t (F^2+1)/(F^2+i).
\end{verbatim}

Every factor is positive and at most one, so F\textless{}\allowbreak{}=\allowbreak{}H by (8). Every factor is nondecreasing in F\textasciicircum{}\allowbreak{}2; since F\textgreater{}\allowbreak{}=\allowbreak{}1,



\begin{verbatim}
f>=2^t/(t+1)>t^2.                             (9)
\end{verbatim}

The strict inequality at t=\allowbreak{}16 is 65536\textgreater{}\allowbreak{}4352. If 2\textasciicircum{}\allowbreak{}t\textgreater{}\allowbreak{}t\textasciicircum{}\allowbreak{}2(t+\allowbreak{}1), then 2\textasciicircum{}\allowbreak{}(t+\allowbreak{}1)\textgreater{}\allowbreak{}2t\textasciicircum{}\allowbreak{}2(t+\allowbreak{}1)\textgreater{}\allowbreak{}=\allowbreak{}(t+\allowbreak{}1)\textasciicircum{}\allowbreak{}2(t+\allowbreak{}2), because t\textasciicircum{}\allowbreak{}2\textgreater{}\allowbreak{}=\allowbreak{}3t+\allowbreak{}2 for t\textgreater{}\allowbreak{}=\allowbreak{}4. Hence (9) holds for all t\textgreater{}\allowbreak{}=\allowbreak{}16. In particular F\textgreater{}\allowbreak{}t\textasciicircum{}\allowbreak{}2.



Let a\_\allowbreak{}i=\allowbreak{}(i-1)/\allowbreak{}(F\textasciicircum{}\allowbreak{}2+\allowbreak{}i). Then



\begin{verbatim}
sum_i a_i < t(t-1)/(2F^2)<1/(2F).
\end{verbatim}

The finite product inequality product\_\allowbreak{}i(1-a\_\allowbreak{}i)\textgreater{}\allowbreak{}=\allowbreak{}1-sum\_\allowbreak{}i a\_\allowbreak{}i gives



\begin{verbatim}
H(1-1/(2F))<f<=F.
\end{verbatim}

Thus



\begin{verbatim}
H<2F^2/(2F-1)=F+F/(2F-1)<=F+1.
\end{verbatim}

Integrality gives H\textless{}\allowbreak{}=\allowbreak{}F. Together with F\textless{}\allowbreak{}=\allowbreak{}H, this proves F=\allowbreak{}t!.



Conversely let F=\allowbreak{}H=\allowbreak{}t!. Since H\textgreater{}\allowbreak{}t\textasciicircum{}\allowbreak{}2, the same estimate gives H-1/\allowbreak{}2\textless{}\allowbreak{}f\textless{}\allowbreak{}=\allowbreak{}H, so (8) and both natural ceiling witnesses exist. Thus the proposed unified graph admits exactly the intended factorial value, not an additional small fixed point.



\subsection*{5. The same B is the reciprocal CRT product}



With F=\allowbreak{}t!, each i in 1..t divides F, and



\begin{verbatim}
d_i=1+F^2/i
\end{verbatim}

is an integer congruent to 1 modulo F. Every prime dividing d\_\allowbreak{}i is greater than t. If a prime p divided two different d\_\allowbreak{}i,d\_\allowbreak{}j, it would divide i\textasteriskcentered{}\allowbreak{}d\_\allowbreak{}i-j\textasteriskcentered{}\allowbreak{}d\_\allowbreak{}j=\allowbreak{}i-j, contradicting p\textgreater{}\allowbreak{}t. The d\_\allowbreak{}i are pairwise coprime and



\begin{verbatim}
product_(i=1)^t d_i=binom(F^2+t,t)=B.
\end{verbatim}

The factorial binomial and the required CRT product are therefore literally the same integer and the same digit graph, not just two equal outputs retained with an extra equality.



\subsection*{6. All other low digits fit the pure radix}



The inherited local majorant certificate has E=\allowbreak{}8388608. The label certificate has ell=\allowbreak{}4194304 for the derived-label variant. Its raw positive-gate ledger gives local L and J less than R\textasciicircum{}\allowbreak{}E. Because F\textgreater{}\allowbreak{}=\allowbreak{}16!\textgreater{}\allowbreak{}2\textasciicircum{}\allowbreak{}44, R\textgreater{}\allowbreak{}=\allowbreak{}2\textasciicircum{}\allowbreak{}44. At r0=\allowbreak{}2\textasciicircum{}\allowbreak{}44, 44E\textless{}\allowbreak{}2\textasciicircum{}\allowbreak{}44, so r0\textasciicircum{}\allowbreak{}E\textless{}\allowbreak{}2\textasciicircum{}\allowbreak{}r0. For every r\textgreater{}\allowbreak{}=\allowbreak{}r0, E/\allowbreak{}r\textless{}\allowbreak{}=\allowbreak{}1/\allowbreak{}2 and



\begin{verbatim}
(1+1/r)^E<=sum_(j=0)^E (E/r)^j<2.
\end{verbatim}

Integer induction gives R\textasciicircum{}\allowbreak{}E\textless{}\allowbreak{}2\textasciicircum{}\allowbreak{}R\textless{}\allowbreak{}2\textasciicircum{}\allowbreak{}K. This is used only after factorial recovery.



For any independent representative n, equation (3) gives n\textless{}\allowbreak{}=\allowbreak{}cN\textasciicircum{}\allowbreak{}2\textless{}\allowbreak{}=\allowbreak{}10N\textasciicircum{}\allowbreak{}2. Every lower range index satisfies 1\textless{}\allowbreak{}=\allowbreak{}k\textless{}\allowbreak{}=\allowbreak{}R-1. For j\textless{}\allowbreak{}=\allowbreak{}k, if j\textgreater{}\allowbreak{}n the coefficient is zero; otherwise



\begin{verbatim}
binom(n,j)<=n^k<=10^k N^(2k)
             <2^(4R) N^(2k)
             <=N^(2k+2)<=N^K<=q.            (10)
\end{verbatim}

The first strict inequality uses k\textless{}\allowbreak{}R and 10\textless{}\allowbreak{}16. The penultimate comparison uses N=\allowbreak{}B\textasciicircum{}\allowbreak{}K\textgreater{}\allowbreak{}=\allowbreak{}2\textasciicircum{}\allowbreak{}K and K=\allowbreak{}2R. Thus every independent range digit is carry-free for q=\allowbreak{}N\textasciicircum{}\allowbreak{}K and a fortiori for q=\allowbreak{}(N\textasciicircum{}\allowbreak{}2)\textasciicircum{}\allowbreak{}K. In both variants N divides q.



For the derived label only, the separate literal ledger bounds J(A)\textless{}\allowbreak{}=\allowbreak{}(n+\allowbreak{}1)\textasciicircum{}\allowbreak{}ell where n=\allowbreak{}max of the four encoded schema parameters. The aggregate gives n\textless{}\allowbreak{}=\allowbreak{}9N\textasciicircum{}\allowbreak{}2 and N\textgreater{}\allowbreak{}=\allowbreak{}10, hence n+\allowbreak{}1\textless{}\allowbreak{}=\allowbreak{}N\textasciicircum{}\allowbreak{}3 and J(A)\textless{}\allowbreak{}=\allowbreak{}N\textasciicircum{}\allowbreak{}(3ell). The exact base inequality 2\textasteriskcentered{}\allowbreak{}16!\textgreater{}\allowbreak{}=\allowbreak{}3ell\textasteriskcentered{}\allowbreak{}17 and induction give 3ell(t+\allowbreak{}1)\textless{}\allowbreak{}=\allowbreak{}2F\textless{}\allowbreak{}K. Thus for every j\textless{}\allowbreak{}=\allowbreak{}t+\allowbreak{}1,



\begin{verbatim}
binom(J(A),j)<=J(A)^(t+1)
                   <=N^[3ell(t+1)]<N^K<=q.   (11)
\end{verbatim}

As before, coefficients with j\textgreater{}\allowbreak{}J(A) are zero. Consequently the additional fused graph represents exactly N|binom(J(A),t+\allowbreak{}1). It does not constrain J(A) to satisfy the independent aggregate bound.



\subsection*{7. CRT decoding and the derived-label range}



For a prime p and positive integers e,k, p\textasciicircum{}\allowbreak{}e|binom(a,k) implies a mod p\textasciicircum{}\allowbreak{}e\textless{}\allowbreak{}k. If a\textless{}\allowbreak{}k this is immediate. Otherwise the valuation formula is a sum of terms each in \{0,1\}. If [a-k+\allowbreak{}1,a] contained no multiple of p\textasciicircum{}\allowbreak{}e, every summand for powers p\textasciicircum{}\allowbreak{}j with j\textgreater{}\allowbreak{}=\allowbreak{}e would vanish, leaving valuation at most e-1, a contradiction.



For row i-1 choose p|d\_\allowbreak{}i and m=\allowbreak{}v\_\allowbreak{}p(d\_\allowbreak{}i). The retained raw index congruence B|C+\allowbreak{}F\textasciicircum{}\allowbreak{}2+\allowbreak{}1 gives C=\allowbreak{}i-1 modulo p. The index range N|binom(C,t) gives C=\allowbreak{}s modulo p\textasciicircum{}\allowbreak{}(Km), 0\textless{}\allowbreak{}=\allowbreak{}s\textless{}\allowbreak{}t. Since p\textgreater{}\allowbreak{}t, necessarily s=\allowbreak{}i-1. Every column range similarly gives A\_\allowbreak{}j=\allowbreak{}w\_\allowbreak{}j modulo p\textasciicircum{}\allowbreak{}(Km) in its intended bounded interval. The nonnegative local polynomial satisfies N|L(C,A), so p\textasciicircum{}\allowbreak{}(Km)|L(i-1,w). Its bound 0\textless{}\allowbreak{}=\allowbreak{}L\textless{}\allowbreak{}2\textasciicircum{}\allowbreak{}K\textless{}\allowbreak{}=\allowbreak{}p\textasciicircum{}\allowbreak{}(Km) forces L=\allowbreak{}0.



For the derived-label variant, N|binom(J(A),t+\allowbreak{}1) gives J(A)=\allowbreak{}s modulo p\textasciicircum{}\allowbreak{}(Km) with 0\textless{}\allowbreak{}=\allowbreak{}s\textless{}\allowbreak{}=\allowbreak{}t. Polynomial congruence gives J(w)=\allowbreak{}s modulo p\textasciicircum{}\allowbreak{}(Km). Since both J(w) and s lie in [0,2\textasciicircum{}\allowbreak{}K), they are equal. Thus every decoded row has J(w)\textless{}\allowbreak{}=\allowbreak{}t without a ninth independent local coordinate. For the local-slack variant, the unchanged local square directly gives t=\allowbreak{}w+\allowbreak{}J.



Conversely choose a bounded tuple for every row and use CRT at precision d\_\allowbreak{}i\textasciicircum{}\allowbreak{}(2K), choosing positive representatives in [1,B\textasciicircum{}\allowbreak{}(2K)]=\allowbreak{}[1,N\textasciicircum{}\allowbreak{}2]. For a range lower index k\textless{}\allowbreak{}=\allowbreak{}K, a falling-factorial numerator has a factor divisible by p\textasciicircum{}\allowbreak{}(2Km); division by k! loses at most k-1\textless{}\allowbreak{}K\textless{}\allowbreak{}=\allowbreak{}Km in p-valuation. This proves all required range divisibilities. The same argument applies to J(A) when J(w)\textless{}\allowbreak{}=\allowbreak{}t, because polynomial congruence gives J(A)=\allowbreak{}J(w) modulo p\textasciicircum{}\allowbreak{}(2Km). The aggregate slack exists since only independent representatives are included. The raw index and local divisibilities follow from the same CRT congruences.



\subsection*{8. Full existence equivalence and source calibration}



Given a natural zero, Sections 3-6 interpret every power, the only exact binomial, the factorial, the product, and every fused range graph in noncircular order. Section 7 yields a valid bounded source row at every i\textless{}\allowbreak{}t with all formula labels \textless{}\allowbreak{}=\allowbreak{}t. The unchanged beta equations identify adjacent positive accumulator states, starting at W\_\allowbreak{}1=\allowbreak{}1+\allowbreak{}F and ending at 1, where W\_\allowbreak{}j=\allowbreak{}1+\allowbreak{}jF. Positive labels \textless{}\allowbreak{}=\allowbreak{}t have pairwise-coprime weights. Multiplication of the selected valid Horn-rule equations and prime valuations gives one net production of falsity and balance at every other nonzero label. If falsity were outside the closure of truth under those rules, counting conclusions outside the closure would be at most the number of outside premises, contradicting that net production. Thus the decoded rules contain a finite C222 refutation. This argument does not require acyclic supplied rule multisets and does not mistake the unit weight of truth for a nonzero label.



Conversely start from a finite calibrated refutation, expand it to a finite proof tree, and choose the padded t\textgreater{}\allowbreak{}=\allowbreak{}16 beyond its length, its finite schema parameters and all complete label sums J, before defining F. Set F=\allowbreak{}t!, construct the pending-multiset positive accumulator path with truth padding, and choose the beta encoding x,y. Sections 5 and 7 provide B,N and all independent representatives. Set V and the chosen pure radix q. Sections 4 and 6 provide the natural ceiling and digit witnesses. The inherited simultaneous Pell theorem supplies all eighteen power graphs at once. Every residual is zero. The exact nonnegative projection restores every eliminated coordinate, preserving both directions.



The C222/\allowbreak{}ZF/\allowbreak{}ZFC calibration is unchanged and explicitly retained from the source archive. Finite coefficient checks and finite semantic regression tests are not substitutes for this unbounded argument and do not decide solvability.



\subsection*{9. Artifacts and reproduction}



Run build\_\allowbreak{}unified\_\allowbreak{}product.py with --output DEST. The default is the derived-label Cantor syntax and q=\allowbreak{}N\textasciicircum{}\allowbreak{}K. Select --syntax xx0 for the square-shell syntax, --labels slack for a ninth local label slack, and --squared-radix-base for q=\allowbreak{}(N\textasciicircum{}\allowbreak{}2)\textasciicircum{}\allowbreak{}K. The endpoint includes raw/\allowbreak{}equation.json, raw/\allowbreak{}raw\_\allowbreak{}audit\_\allowbreak{}input.json, a literal majorant ledger, the natural projection certificate, the whole expanded quartic and an independent exact projection receipt.



The generator imports only the previously reviewed round14 release read-only. Every new artifact is written under the explicit output directory. SOURCE\_\allowbreak{}BINDING.json records the imported executable closure and its bytes; REPRODUCTION.json records byte-exact normal/\allowbreak{}optimized rebuilds. The natural equation is sum(r\_\allowbreak{}i\textasciicircum{}\allowbreak{}2)=\allowbreak{}0 over its complete residual list, not an abbreviated schematic or a square of a sum.

% Complete proof from direct/PROOF.md; literal text, no Markdown math interpretation.
\section{Residue-one orientation and selectively even Pell indices}



\subsection*{0. Fixed equation and scope}



The principal polynomial is the sum of the squares of all 290 literal quadratic residuals in \texttt{integer442/\allowbreak{}equation.json}. Its variable order is the order in that file. It has 442 unrestricted integer variables, exact degree four, 4,652 collected monomials, and coefficient height 128,446. Every collected coefficient is supplied in \texttt{integer442/\allowbreak{}expanded\_\allowbreak{}polynomial.json}. The companion \texttt{integer446/\allowbreak{}} has 446 integer variables, 291 residuals, 4,775 monomials and the same height. These are ordinary fixed integer polynomials, not equations with exponentiation primitives.



The 442-variable source is \texttt{../\allowbreak{}architecture/\allowbreak{}unified\_\allowbreak{}derived\_\allowbreak{}shell/\allowbreak{}raw/\allowbreak{}}; the 446-variable source is \texttt{../\allowbreak{}architecture/\allowbreak{}unified\_\allowbreak{}slack\_\allowbreak{}shell/\allowbreak{}raw/\allowbreak{}}. The former uses a derived label range, the latter an additional local label coordinate. Their ordinary natural-domain arithmetic proof is \texttt{../\allowbreak{}architecture/\allowbreak{}UNIFIED\_\allowbreak{}PRODUCT\_\allowbreak{}PROOF.md}. It identifies the factorial's binomial with the reciprocal prime-power product:



\begin{verbatim}
t = t0+16, rho=F^2, B=binom(F^2+t,t), V0=(F^2+1)^t,
F B=V0+s, B=r+s+1, K=2R, N=B^K, q=N^K.
\end{verbatim}

Here these formulas name intended meanings; powers and binomials are not assumed before their certification. The actual equations have 18 power requests and 15 seed indices. The new compiler keeps the t and K seeds undoubled, doubles the other thirteen, and leaves the derived t+\allowbreak{}1 pair as its original one-step Pell addition. Every one of the eleven \texttt{(q+\allowbreak{}1)\textasciicircum{}\allowbreak{}k} requests is among the doubled seeds. The two other doubled seeds certify q\textasciicircum{}\allowbreak{}(y+\allowbreak{}1) and q\textasciicircum{}\allowbreak{}(x+\allowbreak{}1).



The theorem is an ordinary proof, formalizable in ZFC, that the new polynomial has an integer zero iff the inherited direct source has a natural zero. Composing with that source's numerical C222/\allowbreak{}refutation calibration gives



\begin{verbatim}
ZFC proves: [there is no z in Z^442 with P442(z)=0] iff Con(ZFC).
\end{verbatim}

This is not a proof of consistency, a decision of the equation's solvability, a globally smallest equation, or a newly emitted full Lean/\allowbreak{}object-ZFC proof. The named calibration interfaces remain KEndpoint.c222\_\allowbreak{}iff\_\allowbreak{}zf\_\allowbreak{}refutation and ConstructibleChoice.zf\_\allowbreak{}refutation\_\allowbreak{}iff\_\allowbreak{}zfc\_\allowbreak{}refutation. The arithmetic route does not simulate K215 and does not instantiate Jones's universal polynomial.



\subsection*{1. The two new interface ideas}



\subsubsection*{1.1 Residue one orients a signed square root}



Suppose q\textgreater{}\allowbreak{}=\allowbreak{}3, k\textgreater{}\allowbreak{}=\allowbreak{}1 and m is an integer satisfying



\begin{verbatim}
m^2=(q+1)^(2k),  m=1 (mod q).
\end{verbatim}

The first equality gives m=\allowbreak{}+\allowbreak{}(q+\allowbreak{}1)\textasciicircum{}\allowbreak{}k or m=\allowbreak{}-(q+\allowbreak{}1)\textasciicircum{}\allowbreak{}k. Their residues modulo q are +\allowbreak{}1 and -1. Since q does not divide 2, only the positive sign is possible. Therefore m=\allowbreak{}(q+\allowbreak{}1)\textasciicircum{}\allowbreak{}k. This needs no nonnegativity guard for m.



The first materialized variant adds the single quadratic row m-1-q\textasteriskcentered{}\allowbreak{}u=\allowbreak{}0 for each of eleven outputs, with signed u. The preferred variant builds the same residue condition into an existing remainder. The actual front equation has one of the forms



\begin{verbatim}
m=L*(C+q*h)+r,     or     m=N*L*h+r,
L=r+s+1.
\end{verbatim}

Replace the old remainder r everywhere by 1+\allowbreak{}q\textasteriskcentered{}\allowbreak{}u, reusing its coordinate as u. The equations become



\begin{verbatim}
m=L*(C+q*h)+1+q*u,  or  m=N*L*h+1+q*u,
L=q*u+s+2.
\end{verbatim}

All changed residuals remain quadratic: the old remainder occurs linearly in every source row, and q is an existing coordinate. The certificate lists every occurrence and its exact substitution. Once L is proved to be a positive power of q, the first equation gives m=\allowbreak{}1 modulo q.



This substitution is NOT a bijective change of arbitrary integer coordinates. It is an existence-preserving restriction justified at the semantic interface. For completeness, a genuine terminal power has m=\allowbreak{}(q+\allowbreak{}1)\textasciicircum{}\allowbreak{}k=\allowbreak{}1 modulo q and the old divisor L is q\textasciicircum{}\allowbreak{}j with j\textgreater{}\allowbreak{}=\allowbreak{}1. Thus its old bounded remainder satisfies r=\allowbreak{}1 modulo q, 0\textless{}\allowbreak{}=\allowbreak{}r\textless{}\allowbreak{}L, hence r\textgreater{}\allowbreak{}=\allowbreak{}1; u=\allowbreak{}(r-1)/\allowbreak{}q is a nonnegative integer. For soundness, u\textgreater{}\allowbreak{}=\allowbreak{}0 and q\textgreater{}\allowbreak{}0 give the old nonnegative remainder r=\allowbreak{}1+\allowbreak{}q\textasteriskcentered{}\allowbreak{}u. The recovery of the genuine powers below establishes exactly the required interface.



\subsubsection*{1.2 Weighted interval guards share one sign anchor}



For q\textgreater{}\allowbreak{}0 suppose several pairs satisfy



\begin{verbatim}
q*u_i+s_i=L-2.
\end{verbatim}

Retain u\_\allowbreak{}0\textgreater{}\allowbreak{}=\allowbreak{}0 and s\_\allowbreak{}0\textgreater{}\allowbreak{}=\allowbreak{}0. They imply L-2\textgreater{}\allowbreak{}=\allowbreak{}0. For any other pair,



\begin{verbatim}
[u_i>=0 and s_i>=0] iff u_i*s_i>=0.
\end{verbatim}

Indeed, if either were negative, a nonnegative product would make the other nonpositive, and q\textasteriskcentered{}\allowbreak{}u\_\allowbreak{}i+\allowbreak{}s\_\allowbreak{}i would be strictly negative. This contradicts L-2\textgreater{}\allowbreak{}=\allowbreak{}0. The reverse direction is immediate. This weighted lemma is the exact adaptation of the incoming Structural Compression shared-remainder lemma to the new remainder coordinates. Seven literal applications replace fourteen individual sign predicates by seven product predicates. Every pair, retained anchor, common-bound identity, and row is recorded. The sign of q is established before using this lemma; it does not depend on a binomial interpretation.



All sign predicates G\textgreater{}\allowbreak{}=\allowbreak{}0 are represented by one quadratic row



\begin{verbatim}
G=a^2+b^2+c(c+1),  a,b,c in Z.
\end{verbatim}

This expression takes only nonnegative integer values. Conversely 4G+\allowbreak{}1 is not of the excluded three-square shape; Legendre's three-square theorem and parity put its representation in this form. G is affine or quadratic here, so the final residuals remain quadratic. There are 38 such predicates in the principal source and 39 in the companion, including a strict common-ordinate predicate.



\subsection*{2. Pell notation and facts used in the compiler}



For a\textgreater{}\allowbreak{}=\allowbreak{}2 set D=\allowbreak{}a\textasciicircum{}\allowbreak{}2-1 and define



\begin{verbatim}
(a+sqrt(D))^j=chi_a(j)+psi_a(j)*sqrt(D),  j>=0.
\end{verbatim}

They are integer sequences satisfying the usual multiplication recurrence, with (chi(0),psi(0))=\allowbreak{}(1,0) and (chi(1),psi(1))=\allowbreak{}(a,1). Both are strictly increasing after the zero endpoint, and psi\_\allowbreak{}a(j)\textgreater{}\allowbreak{}=\allowbreak{}j. Every integral solution of x\textasciicircum{}\allowbreak{}2-D\textasteriskcentered{}\allowbreak{}y\textasciicircum{}\allowbreak{}2=\allowbreak{}1 has |x|=\allowbreak{}chi\_\allowbreak{}a(i), |y|=\allowbreak{}psi\_\allowbreak{}a(i) for one i\textgreater{}\allowbreak{}=\allowbreak{}0. To prove this, reduce a positive norm-one unit by powers of a+\allowbreak{}sqrt(D); the smallest nontrivial positive integral ordinate is 1 and its abscissa is a. The standard reduction therefore leaves only the identity. Independent signs of the coordinates are allowed.



Two elementary facts supply the finite indexed-Pell relation.



\begin{enumerate}

\item If i,h\textgreater{}\allowbreak{}=\allowbreak{}1 and psi\_\allowbreak{}a(i)\textasciicircum{}\allowbreak{}2 divides psi\_\allowbreak{}a(h), then i divides h and psi\_\allowbreak{}a(i) divides h. Write h=\allowbreak{}li+\allowbreak{}r with 0\textless{}\allowbreak{}=\allowbreak{}r\textless{}\allowbreak{}i, x=\allowbreak{}chi\_\allowbreak{}a(i), y=\allowbreak{}psi\_\allowbreak{}a(i). Modulo y the addition law gives psi\_\allowbreak{}a(h)=\allowbreak{}x\textasciicircum{}\allowbreak{}l psi\_\allowbreak{}a(r). The norm gives gcd(x,y)=\allowbreak{}1; monotonicity forces r=\allowbreak{}0. Expansion modulo y\textasciicircum{}\allowbreak{}2 then gives psi\_\allowbreak{}a(li)=\allowbreak{}l\textasteriskcentered{}\allowbreak{}x\textasciicircum{}\allowbreak{}(l-1)\textasteriskcentered{}\allowbreak{}y, so y divides l, hence y divides h. Also i\textless{}\allowbreak{}=\allowbreak{}h.

\item For 1\textless{}\allowbreak{}=\allowbreak{}i\textless{}\allowbreak{}=\allowbreak{}h, if chi\_\allowbreak{}a(j)=\allowbreak{}+\allowbreak{}/\allowbreak{}-chi\_\allowbreak{}a(i) modulo chi\_\allowbreak{}a(h), then j=\allowbreak{}+\allowbreak{}/\allowbreak{}-i modulo 2h. Modulo u=\allowbreak{}chi\_\allowbreak{}a(h), the norm-unit doubling identity gives (a+\allowbreak{}sqrt(D))\textasciicircum{}\allowbreak{}(2h)=\allowbreak{}-1. Reduce j modulo 2h and reflect its remainder about h. This leaves +\allowbreak{}/\allowbreak{}-chi\_\allowbreak{}a(r) with 0\textless{}\allowbreak{}=\allowbreak{}r\textless{}\allowbreak{}=\allowbreak{}h and r=\allowbreak{}+\allowbreak{}/\allowbreak{}-j modulo 2h. Since 2\textasteriskcentered{}\allowbreak{}chi\_\allowbreak{}a(h-1)\textless{}\allowbreak{}chi\_\allowbreak{}a(h), two coordinates below h cannot add to u. Congruence up to sign forces r=\allowbreak{}i. The endpoint i=\allowbreak{}h is the unique zero residue. In particular congruence to +\allowbreak{}/\allowbreak{}-a forces j odd.

\end{enumerate}



These statements are unbounded arithmetic lemmas; finite Pell tests do not replace their proofs. Polynomial dependence of chi\_\allowbreak{}b(j),psi\_\allowbreak{}b(j) on b permits base transfer modulo any integer. At b=\allowbreak{}1 one has chi\_\allowbreak{}1(j)=\allowbreak{}1, psi\_\allowbreak{}1(j)=\allowbreak{}j.



\subsection*{3. A common square frame and common nonzero ordinate}



Let Aset be the distinct actual base, exponent, and output atoms of the eighteen original power requests, using their exact affine restorations in the contracted source. Replace the private frame by



\begin{verbatim}
W=2*(2+sum(a^2 for a in Aset)).
\end{verbatim}

Consequently W\textgreater{}\allowbreak{}=\allowbreak{}4, |a|\textless{}\allowbreak{}W, a\textasciicircum{}\allowbreak{}2\textless{}\allowbreak{}W, and 2|a|\textless{}\allowbreak{}W for each atom. In particular the requested doubled exponents 2k and squared outputs m\textasciicircum{}\allowbreak{}2 are strictly below W. No sign of a proposed power output is needed for this frame. Its exact affine atom list and defining row are in the transformation certificate.



Retain A\textgreater{}\allowbreak{}=\allowbreak{}0 and Zminus1\textgreater{}\allowbreak{}=\allowbreak{}0 and the old frame norm



\begin{verbatim}
A^2-W^3(W+2)*(Zminus1+1)^2=1.
\end{verbatim}

Set Z=\allowbreak{}Zminus1+\allowbreak{}1\textgreater{}\allowbreak{}=\allowbreak{}1. Pell classification at center W+\allowbreak{}1 identifies (A,WZ)=\allowbreak{}(chi\_\allowbreak{}(W+\allowbreak{}1)(h),psi\_\allowbreak{}(W+\allowbreak{}1)(h)) with h\textgreater{}\allowbreak{}=\allowbreak{}1. Reduction modulo W gives psi\_\allowbreak{}(W+\allowbreak{}1)(h)=\allowbreak{}h modulo W, so W divides h and h\textgreater{}\allowbreak{}=\allowbreak{}W. Therefore



\begin{verbatim}
A >= (W+1)^W > W^W,  A>W^2.
\end{verbatim}

For every positive actual power base 1\textless{}\allowbreak{}=\allowbreak{}n\textless{}\allowbreak{}W, the existing definition



\begin{verbatim}
T_n=2An-n^2-1
\end{verbatim}

gives T\_\allowbreak{}n\textgreater{}\allowbreak{}=\allowbreak{}2A-2\textgreater{}\allowbreak{}A\textgreater{}\allowbreak{}W\textasciicircum{}\allowbreak{}W. This follows because the quadratic increases on 1\textless{}\allowbreak{}=\allowbreak{}n\textless{}\allowbreak{}W\textless{}\allowbreak{}A and its minimum is at n=\allowbreak{}1.



Keep the common norm U\textasciicircum{}\allowbreak{}2-(A\textasciicircum{}\allowbreak{}2-1)Vaux\textasciicircum{}\allowbreak{}2=\allowbreak{}1, the original product identity



\begin{verbatim}
Vaux=(product over fifteen seeds Y_i^2)*(v+1),
\end{verbatim}

and impose Vaux\textgreater{}\allowbreak{}=\allowbreak{}1 directly. Delete the old sign predicate v\textgreater{}\allowbreak{}=\allowbreak{}0. The product and positivity imply all Y\_\allowbreak{}i are nonzero, their squares divide Vaux, and v+\allowbreak{}1\textgreater{}\allowbreak{}0, hence v\textgreater{}\allowbreak{}=\allowbreak{}0. Thus this change supplies its own omitted sign before using the index theorem. The common norm implies |U|\textgreater{}\allowbreak{}=\allowbreak{}A\textgreater{}\allowbreak{}=\allowbreak{}2.



Retain the constructor Baux=\allowbreak{}1+\allowbreak{}4Vaux\textasteriskcentered{}\allowbreak{}(q4+\allowbreak{}1) with q4\textgreater{}\allowbreak{}=\allowbreak{}0 and Baux=\allowbreak{}A modulo U. Thus Baux\textgreater{}\allowbreak{}=\allowbreak{}5 is odd, and Baux=\allowbreak{}1 modulo 4|Y\_\allowbreak{}i| for every seed. No auxiliary comparison coordinate or congruence quotient needs a sign assumption.



\subsection*{4. Even seed indices recover their own missing size bounds}



For a selected seed, k\textgreater{}\allowbreak{}=\allowbreak{}1 is its positive original exponent and e=\allowbreak{}2k\textless{}\allowbreak{}W. Keep its three norms and abscissa congruence, but change its ordinate congruence to



\begin{verbatim}
X^2-(A^2-1)Y^2=1,
S^2-(Baux^2-1)Raux^2=1,
S=X (mod U),  Raux=2k (mod 4Y).
\end{verbatim}

The private old row Y=\allowbreak{}k+\allowbreak{}slack may be removed with its unconstrained slack. It has no other occurrence after its sign guard is removed. Its unique inverse is slack=\allowbreak{}Y-k; this exact pruning is separately certified.



We already know Y!=\allowbreak{}0. Suppose |Y|=\allowbreak{}1. Then |X|=\allowbreak{}A. Classify the comparison norm by S=\allowbreak{}+\allowbreak{}/\allowbreak{}-chi\_\allowbreak{}Baux(j), Raux=\allowbreak{}+\allowbreak{}/\allowbreak{}-psi\_\allowbreak{}Baux(j). Transfer Baux=\allowbreak{}A modulo U and S=\allowbreak{}X modulo U. The abscissa residue fact implies that j is odd. Since Baux is odd, psi\_\allowbreak{}Baux(j)=\allowbreak{}j modulo 2; hence Raux is odd. But Raux=\allowbreak{}2k modulo 4Y makes it even. Contradiction. Therefore |Y|\textgreater{}\allowbreak{}=\allowbreak{}2. Its main Pell index i is at least two, and



\begin{verbatim}
|Y|=psi_A(i)>=psi_A(2)=2A>2k.
\end{verbatim}

Now classify the common norm as |U|=\allowbreak{}chi\_\allowbreak{}A(h), Vaux=\allowbreak{}psi\_\allowbreak{}A(h), h\textgreater{}\allowbreak{}=\allowbreak{}1. Squared divisibility gives i\textless{}\allowbreak{}=\allowbreak{}h and |Y| divides h. The S congruence yields j=\allowbreak{}+\allowbreak{}/\allowbreak{}-i modulo 2h. Since Baux=\allowbreak{}1 modulo 4|Y|, the Raux congruence yields j=\allowbreak{}+\allowbreak{}/\allowbreak{}-2k modulo 4|Y|. Therefore i=\allowbreak{}+\allowbreak{}/\allowbreak{}-2k modulo 2|Y|. But 1\textless{}\allowbreak{}=\allowbreak{}i\textless{}\allowbreak{}=\allowbreak{}|Y| and 1\textless{}\allowbreak{}=\allowbreak{}2k\textless{}\allowbreak{}|Y|; the minus congruence is impossible and the plus congruence forces i=\allowbreak{}2k. Thus



\begin{verbatim}
|X|=chi_A(2k), |Y|=psi_A(2k).
\end{verbatim}

For the two undoubled t and K seeds, their retained guards Y\textgreater{}\allowbreak{}=\allowbreak{}k\textgreater{}\allowbreak{}=\allowbreak{}1 give the usual same short-interval proof, with positive Y and |X|=\allowbreak{}chi\_\allowbreak{}A(k). All exponent positivity is syntactic before power semantics: t\textgreater{}\allowbreak{}=\allowbreak{}16, K=\allowbreak{}2R\textgreater{}\allowbreak{}0, t+\allowbreak{}F\textasciicircum{}\allowbreak{}2\textgreater{}\allowbreak{}0, x+\allowbreak{}1\textgreater{}\allowbreak{}0, y+\allowbreak{}1\textgreater{}\allowbreak{}0, and every encoded representative plus one is positive. The derived label J, when present, has a monotone constructor with positive constant term. Its source certificate supplies the compositional lower bound.



\subsection*{5. Squared outputs force positive Pell coordinates}



For a doubled seed and a supporting base n\textgreater{}\allowbreak{}=\allowbreak{}2, replace its power congruence by



\begin{verbatim}
X+(n-A)Y=m^2+T_n*h.
\end{verbatim}

Set e=\allowbreak{}2k, N0=\allowbreak{}n\textasciicircum{}\allowbreak{}e, M=\allowbreak{}m\textasciicircum{}\allowbreak{}2, x=\allowbreak{}chi\_\allowbreak{}A(e), y=\allowbreak{}psi\_\allowbreak{}A(e). The standard recurrence and norm give



\begin{verbatim}
H=x+(n-A)y = N0 (mod T_n),
J=x-(n-A)y,  HJ=1+T_n*y^2.
\end{verbatim}

The four independent signs of X,Y imply respectively that T\_\allowbreak{}n divides



\begin{verbatim}
M-N0,  M+N0,  M*N0-1,  M*N0+1.
\end{verbatim}

The frame has n,e,M\textless{}\allowbreak{}W. Hence N0\textless{}\allowbreak{}W\textasciicircum{}\allowbreak{}(W-1), M\textasteriskcentered{}\allowbreak{}N0\textless{}\allowbreak{}W\textasciicircum{}\allowbreak{}W\textless{}\allowbreak{}A\textless{}\allowbreak{}T\_\allowbreak{}n, and M+\allowbreak{}N0\textless{}\allowbreak{}W+\allowbreak{}W\textasciicircum{}\allowbreak{}(W-1)\textless{}\allowbreak{}W\textasciicircum{}\allowbreak{}W\textless{}\allowbreak{}T\_\allowbreak{}n. The second and fourth divisibilities are impossible by strict positivity. The third is impossible because N0\textgreater{}\allowbreak{}=\allowbreak{}4 and M is a nonnegative integer, so M\textasteriskcentered{}\allowbreak{}N0 cannot equal one. Only X=\allowbreak{}x,Y=\allowbreak{}y remains, and the first congruence gives M=\allowbreak{}N0 exactly. Thus m=\allowbreak{}+\allowbreak{}/\allowbreak{}-n\textasciicircum{}\allowbreak{}k and the actual main Pell coordinates are positive, even before the sign of m is resolved.



For q-power outputs already known positive from their remainder bounds, this immediately gives the intended power. For terminal successor-base outputs, Section 1.1 gives the positive sign once their front division yields m=\allowbreak{}1 modulo q. No factorial or binomial value is assumed to obtain this orientation.



\subsection*{6. The noncircular soundness order}



The exact sign-inference certificate uses only integer polynomial identities, positive monomials, and identified affine squares. Its omitted paired predicates are restored first by the weighted interval lemma. The needed q\textgreater{}\allowbreak{}0 is independent: the unchanged factorial remainder/\allowbreak{}slack predicates give



\begin{verbatim}
B=fact_remainder+fact_slack+1>=1,
q=B+choose_oslack+1>=2.
\end{verbatim}

The retained anchor for the binomial remainder gives



\begin{verbatim}
L=q*u+choose_rslack+2>=2.
\end{verbatim}

Every other interval anchor is positive in the same way. No power is interpreted in these statements. After all seven pair restorations, the emitted algebraic closure has every required ordinary front sign. Its selected Y-slacks, comparison coordinates, ordinary division quotients, and power outputs need no such sign.



\begin{enumerate}

\item The retained t-seed guard identifies Y=\allowbreak{}psi\_\allowbreak{}A(t)\textgreater{}\allowbreak{}0 and X=\allowbreak{}+\allowbreak{}/\allowbreak{}-chi\_\allowbreak{}A(t). Its request at base q has output L\textgreater{}\allowbreak{}=\allowbreak{}2. The signed power residue excludes the negative X: multiplying by q\textasciicircum{}\allowbreak{}t would force L\textasteriskcentered{}\allowbreak{}q\textasciicircum{}\allowbreak{}t=\allowbreak{}-1 under the strict frame bound, impossible. Hence L=\allowbreak{}q\textasciicircum{}\allowbreak{}t and this t pair is positively oriented.

\item The other t request now gives V0=\allowbreak{}(F\textasciicircum{}\allowbreak{}2+\allowbreak{}1)\textasciicircum{}\allowbreak{}t\textgreater{}\allowbreak{}0, even if the base was initially one. The literal ceiling F\textasteriskcentered{}\allowbreak{}B=\allowbreak{}V0+\allowbreak{}fact\_\allowbreak{}slack and B\textgreater{}\allowbreak{}=\allowbreak{}1 imply F\textgreater{}\allowbreak{}=\allowbreak{}1. Moreover V0=\allowbreak{}(F\textasciicircum{}\allowbreak{}2+\allowbreak{}1)\textasciicircum{}\allowbreak{}t\textgreater{}\allowbreak{}=\allowbreak{}F\textasciicircum{}\allowbreak{}2+\allowbreak{}1, so F\textasteriskcentered{}\allowbreak{}B\textgreater{}\allowbreak{}=\allowbreak{}F\textasciicircum{}\allowbreak{}2+\allowbreak{}1 and B\textgreater{}\allowbreak{}F\textgreater{}\allowbreak{}=\allowbreak{}1. Thus B\textgreater{}\allowbreak{}=\allowbreak{}2 and q\textgreater{}\allowbreak{}=\allowbreak{}3. This precedes interpreting B as a binomial.

\item The retained K-seed guard and request N=\allowbreak{}B\textasciicircum{}\allowbreak{}K at base B\textgreater{}\allowbreak{}=\allowbreak{}2 orient that pair by the ordinary signed-output small-residue argument. Its other request is then q=\allowbreak{}N\textasciicircum{}\allowbreak{}K at positive base N\textgreater{}\allowbreak{}=\allowbreak{}2. Both equalities are now genuine.

\item The derived t+\allowbreak{}1 pair is still the original one-step addition from the positively oriented t pair. Its q\textasciicircum{}\allowbreak{}(t+\allowbreak{}1) output is correct. The two selected q-power seeds use Section 5 and their positive interval outputs to recover q\textasciicircum{}\allowbreak{}(y+\allowbreak{}1) and q\textasciicircum{}\allowbreak{}(x+\allowbreak{}1). All front divisors L are therefore positive powers of q of exponent at least one, and hence divisible by q.

\item Apply Sections 4-5 to the remaining eleven selected seeds. Their outputs have |m|=\allowbreak{}(q+\allowbreak{}1)\textasciicircum{}\allowbreak{}k. Every front digit/\allowbreak{}division equation gives m=\allowbreak{}1 modulo q through the explicit substituted remainder. Since q\textgreater{}\allowbreak{}=\allowbreak{}3, the negative sign is impossible. All eighteen original power equalities have been recovered.

\item Restore each original front remainder r=\allowbreak{}1+\allowbreak{}q\textasteriskcentered{}\allowbreak{}u\textgreater{}\allowbreak{}=\allowbreak{}1. Its exact old bound is r+\allowbreak{}s+\allowbreak{}1=\allowbreak{}L, so it is nonnegative and bounded. The original high-digit and range quotients, though initially signed, are nonnegative whenever required: m=\allowbreak{}L(C+\allowbreak{}q\textasteriskcentered{}\allowbreak{}h)+\allowbreak{}r, 0\textless{}\allowbreak{}=\allowbreak{}C\textless{}\allowbreak{}q, 0\textless{}\allowbreak{}=\allowbreak{}r\textless{}\allowbreak{}L and m\textgreater{}\allowbreak{}0 exclude h\textless{}\allowbreak{}=\allowbreak{}-1; similarly m=\allowbreak{}N\textasteriskcentered{}\allowbreak{}L\textasteriskcentered{}\allowbreak{}h+\allowbreak{}r with N\textgreater{}\allowbreak{}=\allowbreak{}1 excludes h\textless{}\allowbreak{}=\allowbreak{}-1. All remaining front meanings follow in the source's factorial, exact-binomial, prime-power lifting, CRT-range, and direct-calculus order.

\end{enumerate}



For reference, the ordinary signed-output exclusion used in steps 1 and 3 needs no sign of the candidate m. With positive Y=\allowbreak{}psi\_\allowbreak{}A(k), negative X would give m\textasteriskcentered{}\allowbreak{}n\textasciicircum{}\allowbreak{}k=\allowbreak{}-1 after the small-residue argument. This is impossible for n\textgreater{}\allowbreak{}=\allowbreak{}2. The frame bounds |m|,n,k\textless{}\allowbreak{}W ensure |m\textasteriskcentered{}\allowbreak{}n\textasciicircum{}\allowbreak{}k+\allowbreak{}1|\textless{}\allowbreak{}T\_\allowbreak{}n. Positive X then gives m=\allowbreak{}n\textasciicircum{}\allowbreak{}k. The exceptional base n=\allowbreak{}1 is used only after the pair is already oriented.



This recovers a genuine semantic front tuple. If the source theorem requires its old private Pell witnesses rather than just the semantic interface, rebuild fresh ones at its original indices and frame by Section 7. No pointwise claim that doubled Pell coordinates solve an undoubled system is made.



\subsection*{7. Complete witness construction in the reverse direction}



Begin with a genuine semantic front tuple supplied by the inherited natural source, not with a hypothetical standard-model consistency assumption. Keep its logical/\allowbreak{}arithmetic coordinates, and perform the following finite choices.



\begin{enumerate}

\item For each terminal successor-base division, take its genuine old remainder r and set u=\allowbreak{}(r-1)/\allowbreak{}q. As proved in Section 1.1, this is a nonnegative integer. Keep its old slack s. The new front equations and weighted sign predicates hold. Choose the frame W from the original base/\allowbreak{}exponent/\allowbreak{}output atoms.

\item Put A=\allowbreak{}chi\_\allowbreak{}(W+\allowbreak{}1)(W), Z=\allowbreak{}psi\_\allowbreak{}(W+\allowbreak{}1)(W)/\allowbreak{}W. The recurrence modulo W proves that Z is a positive integer. Set Zminus1=\allowbreak{}Z-1. This supplies the common frame norm and all positive moduli T\_\allowbreak{}n.

\item For each of the two undoubled seeds choose (X,Y)=\allowbreak{}(chi\_\allowbreak{}A(k),psi\_\allowbreak{}A(k)); for each of thirteen selected seeds choose (chi\_\allowbreak{}A(2k),psi\_\allowbreak{}A(2k)). Derive the t+\allowbreak{}1 pair by the unchanged one-step recurrence. Let P be the product of all fifteen seed ordinates squared.

\item In the finite ring (Z/\allowbreak{}(4P)Z)[s]/\allowbreak{}(s\textasciicircum{}\allowbreak{}2-(A\textasciicircum{}\allowbreak{}2-1)), A+\allowbreak{}s is invertible, with inverse A-s. A positive power is one. Its integer Pell coordinates give U=\allowbreak{}1 modulo 4P, Vaux=\allowbreak{}0 modulo 4P, U\textgreater{}\allowbreak{}0, Vaux\textgreater{}\allowbreak{}0. Thus P divides Vaux, and v=\allowbreak{}Vaux/\allowbreak{}P-1\textgreater{}\allowbreak{}=\allowbreak{}0. The Pell norm gives gcd(U,Vaux)=\allowbreak{}1; U is odd, so gcd(U,4Vaux)=\allowbreak{}1.

\item CRT supplies Baux=\allowbreak{}A modulo U, Baux=\allowbreak{}1 modulo 4Vaux. Increase it along that progression until Baux\textgreater{}\allowbreak{}A and Baux\textgreater{}\allowbreak{}1+\allowbreak{}4Vaux. Its two exact quotient witnesses and q4 are then nonnegative. For each actual requested index e, set S=\allowbreak{}chi\_\allowbreak{}Baux(e), Raux=\allowbreak{}psi\_\allowbreak{}Baux(e). Polynomial base transfer supplies S=\allowbreak{}X modulo U and Raux=\allowbreak{}e modulo 4Y. Their quotient witnesses are integers.

\item Every genuine output makes its unchanged or squared power congruence exact. For e=\allowbreak{}k or 2k the identity

\end{enumerate}



\begin{verbatim}
   chi_A(e)+(n-A)psi_A(e)-n^e
     =T_n*sum_{j=0}^{e-1} n^(e-1-j)*psi_A(j)
\end{verbatim}

supplies the integer quotient directly. Fill every ordinary defining gate by its displayed polynomial. The retained t/\allowbreak{}K index slacks are nonnegative. Fill all retained sign triples using the three-square representation.



The choice order is front data, W, A and the main pairs, the common U,Vaux, Baux, then comparison pairs and quotients. No later-chosen modulus is already constrained by a residue depending on its future value. The common finite-ring modulus contains all fifteen ordinate squares, so these choices are simultaneous. This proves completeness. Using undoubled indices and the old frame gives the fresh old certificate needed in the soundness map.



\subsection*{8. Exact algebra after the semantic change}



The guard expansion and private-slack pruning leave a quadratic integer system. The old Horn product is then factored exactly:



\begin{verbatim}
(product L_i^2)*(Bguard^2+G^2)
  =(product L_i)^2*(Bguard^2+G^2).
\end{verbatim}

The signed product is squared before being added to the two independent beta squares; no cancellation argument is used. The guarded factor ((c-a)\textasciicircum{}\allowbreak{}2\textasteriskcentered{}\allowbreak{}(c-b)\textasciicircum{}\allowbreak{}2-(d+\allowbreak{}1)\textasciicircum{}\allowbreak{}2)\textasciicircum{}\allowbreak{}2 is likewise computed by first forming the signed product (c-a)(c-b). The existing round14 factorization certificate records all changed definitions and verifies that no altered intermediate escapes its subnetwork. Its proof applies unchanged to the new front source.



Next, monic integral projection deletes only actual rows z-Q=\allowbreak{}0 (up to sign), with integral Q of degree at most two. Every substitution is made everywhere and checked not to increase any residual degree above two. A quadratic Q is substituted only where z has no nonlinear occurrence outside its defining row. Reverse substitution restores every deleted coordinate; unused coordinates may be restored as zero. The separately pruned Y-slacks restore as Y-k.



Finally the final variables undergo integral translations. In particular old t0=\allowbreak{}new\_\allowbreak{}t-16, not new\_\allowbreak{}t-1, absorbs the source's new lower-bound offset. The additional recorded finite centering reduces coefficient height and support. Every final shift is explicitly listed; old=\allowbreak{}new+\allowbreak{}shift and new=\allowbreak{}old-shift are inverse maps on all integer tuples. These translations, unlike the remainder restriction in Section 1, are genuine integral bijections. A deterministic finite coordinate-descent search found them; no optimality is claimed.



The final polynomial is the sum of the squares of the resulting residuals. An integer tuple makes that sum zero iff every residual is zero. The whole collected expansion is regenerated exactly. The highest occurring degree is four, and every retained variable occurs in the supplied polynomial.



\subsection*{9. Reproduction, dependency boundary, and credit}



Run the following from the round15 package root:



\begin{verbatim}
python direct/build_residue.py architecture/unified_derived_shell/raw direct/rebuilt_source --reparam
python direct/center_integer.py direct/rebuilt_source direct/rebuilt_integer442
\end{verbatim}

Only Python's standard library is required. Inputs are literal source JSON; no upstream historical builder, network fetch, theorem prover, or solution search is executed. The compiler records source hashes, all selected-index changes, frame atoms, complete remainder substitutions, guard predicates, weighted-pair anchors, exact private-slack inverses, Horn factorization, monic projection, and final translations. The polynomial coefficients are reconstructible from the printed residual specification and from these files.



The even-index parity argument and shared Pell completeness are credited to the incoming Carriers and Even Indices package. Ordinary signed-output recovery, square budgets, Horn factorization and exact sign closure are inherited from round14 and related signed-power work. Shared remainder predicates are credited to incoming Structural Compression; the new weighted version accompanies the new quotient remainder. The unified factorial/\allowbreak{}product architecture has its own attribution and proof in the architecture packet. The new contribution here is their selective composition with residue-one orientation, which avoids a new nonnegativity triple for each signed terminal power, and its complete smaller integer quartics. Novelty and comparisons are relative to inspected sources.



Finite checks certify literal algebra and source binding. The unbounded Pell, integer sign, remainder, factorial, CRT, and refutation arguments prove the semantic equivalence. The named inherited formal calibration is retained as a dependency; its full transitive theorem-prover closure was not rebuilt here.

\section{Review of the eight incoming packages}
This review is pinned to the public snapshot named above. The exact archive names, Git blob SHA-1 identifiers, SHA-256 hashes, byte lengths, and acquisition receipts are retained in \texttt{review/SOURCE\_PINS.json} and the adjoining receipt files. No incoming author claim is silently promoted to a newly kernel-checked theorem.

\subsection{Small variable equations}
The following table records degree followed by live operation count. The word positive means strictly positive, not nonnegative.
\begin{center}\small
\begin{tabular}{@{}llp{82mm}@{}}\toprule
Packet&Domain&Reported degree/cost points\\\midrule
Coprime Projection&$11\Z$&688/126, 728/122, 1632/114, 1892/114\\
Coprime Projection&$12\Z$&500/148\\
Coprime Projection&$13\Z,14\Z,15\Z$&274/129, 108/119, 56/110 respectively\\
Powered Saturation round 14&$12\Z$&768/141\\
Saturated Pell&$12\Z$&1084/123; dominated control 1280/123\\
Carriers and Even Indices&$12\Z$&1040/140, 816/144, 800/147, 784/154, 768/161\\
Structural Compression&$13\Z$&536/126\\
Structural Compression&$12\Z$&1200/138, 1456/137\\
Incoming positive frontier&ten positive&283/92, 278/93, 266/96, 258/98, 254/99\\\bottomrule
\end{tabular}\end{center}
The $11\Z$ incoming points already improve the delivered round-14 $12\Z$ point. That fact motivated the present composition rather than another twelve-variable result. The lower-operation $1632/114$ point remains a meaningful tradeoff against $464/119$. The $12\Z$ $500/148$ point is now dominated by the centered eleven-variable result in these measures. The degree-$278$/cost-$93$ positive point appears in overlapping incoming material; matching statistics alone do not establish independent novelty. The ten-positive $266/96$ point improves both the active $360/98$ point and the structural $360/96$ point, but is not an unrestricted-integer equation.

\subsection{Direct quartic families}
Every row below has degree four. The metrics are variables, residuals, collected monomials, and coefficient height.
\begin{center}\small
\begin{tabular}{@{}lrrrr@{}}\toprule
Family and domain&Variables&Rows&Terms&Height\\\midrule
Round 14, $\Z$&519&321&6253&411060\\
Round 14 companion, $\Z$&517&310&6774&3904267\\
Round 14, $\N$&349&268&12768&86915503\\
Signed Powers quotient, $\Z$&577&363&7326&2445790974704\\
Signed Powers low height, $\Z$&582&368&6069&964213\\
Signed square output, $\Z$&586&372&7059&2445790974704\\
Signed square low height, $\Z$&591&377&5835&964213\\
Private units, $\Z$&3564&3350&14518&2\\
Structural, $\Z$&589&371&8156&2445790974704\\
Structural, $\Z$&595&377&7595&6956012\\
Structural, $\Z$&600&382&7075&964213\\
Carriers, $\Z$&566&372&6431&2445790974704\\
Carriers, $\Z$&590&396&4776&964219\\
Carriers, $\Z$&597&403&4767&938267\\
Quadratic Frames, $\Z$&576&368&7172&2445790974704\\
Quadratic Frames, $\Z$&582&374&6200&964213\\\bottomrule
\end{tabular}\end{center}
The Synchronized Exponents packet also gives natural quartics at $358/275/8229$ and $361/278/6759$ in variables/rows/terms, with respective heights
$9903520314287616167992362844$ and $9903520314287616167799424614$. The new $329\N$ endpoint has fewer variables and much smaller coefficients, but more terms than either of those support-focused variants. It is not globally dominant across all metrics. The height-two construction has coefficient alphabet $\{-2,1,2\}$ and constant coefficient one; it addresses a distinct objective.

\subsection{What was independently checked}
The incoming review used a new read-only checker, importing no incoming executable code, to expand every squared residual in ten quartics from the four newly acquired packages. Every coefficient agreed with its supplied polynomial, every advertised variable was live, and exact degree, support, and height matched. It also recomputed topology, closure, and live operation counts for all four eleven-integer Coprime Projection graphs. Their degree claims were assessed through the stated leading-form arguments rather than represented as freshly executed upstream replay. Exact receipts are retained.

The focused arithmetic review of the coprime-fusion proof found its essential interfaces explicit: positive coprime square factors, a carrier with no rational zero, separation of dependent squareclasses, signed folding before projection, even strong multipliers handled honestly, and the corrected period modulo $f^2$. The direct incoming material likewise required careful sign recovery and factorial bootstrap order. The present report and independent audits supply the new compositions, rather than assuming compatibility from matching endpoint metrics.

\subsection{Known failures do not change the retained boundary}
The public H10 README refutes the shared-projection83 proposal on inherited modified75 slices, not the modified75 compiler itself. It separately refutes weakened-bound86 and explicitly preserves established75/certificate87. Those failed candidates omit a necessary projection or a strict positivity bound. They neither prove nor refute the distinct source34 ordinary calibration imported here. The present scalar proof restores the signed source interface before using that calibration, and the direct proof explicitly recovers each required sign and bound before decoding.

The active public direct page records $569$ integer variables, $363$ residuals, $7197$ terms, and height $2445790974704$, plus lower-height alternatives. These supersede stale historical descriptions but did not supersede the separately delivered round-14 $517/519$-variable constructions. The eight ZIP additions are incoming research artifacts, not an automatic replacement of the active page.

\section{Reproduction and the evidence boundary}
The archive is self-contained for finite verification. Its top-level \texttt{verify.py} verifies the file manifest, then copies the packet into a fresh temporary directory. All mutating builders and checkers run there. The exact equation files and complete collected quartics are regenerated from the bundled sources and compared byte for byte. The ordinary proof text is not replaced by sampled tests.

\subsection{Required software and commands}
The release replay passes $31$ commands, including both ordinary and optimized checking modes, and $198$ byte-exact reconstructed-file comparisons; all six coefficient recipes are checked. The finite checks require Python 3.10 or later and SymPy. Standard-library code handles the direct construction, coefficient recipe, hashing, gzip payloads, sparse coefficient expansion, and certificate replay. SymPy is used for independent universal identities and exact scalar register cuts. No network request, original workspace, ProveIt checkout, or proof-assistant installation is needed. The supplied requirements file records the dependency; the verifier never installs software.
\begin{verbatim}
python verify.py --report replay.json
python -O verify.py --report replay-optimized-driver.json
\end{verbatim}
By default each command runs both normal and optimized independent checks. The options \texttt{--mode normal} and \texttt{--mode optimized} select one checking mode. The inherited source34 recipe intentionally refuses optimized Python: it is run in its supported checked mode once, and its complete six coefficient outputs are independently authenticated in both checking modes. This limitation is disclosed rather than describing an unsupported recipe mode as passing.

The report can be built from \texttt{report/main.tex} with pdfLaTeX. The build script uses three passes and packages only the requested PDF and editable TeX sources, not transient auxiliary files or page images. In environments with missing TeX filename databases, pointing \texttt{TEXMF} to the installed TeX Live and Latin Modern trees and initializing a private pdfLaTeX format avoids system installation or mutation.

\subsection{What the finite checks establish}
The scalar audit transcribes the mathematical registers independently, binds all six authentic integer payloads, checks the full universal norm polynomial and exact centering identity, and obtains both a global degree upper bound and a nonzero authentic-specialization lower bound. The literal graph audit checks topology, live inputs, exact operation counts, and the nonzero full leading form. Altered signs, scalings, registers, and final operations must be rejected.

The direct audit imports no generator or eliminator. It verifies the full square frame, actual power-atom expansions, exact seed selection, all remainder substitutions, weighted interval anchors, sign-inference equalities, three-square guards, private-slack restorations, Horn factorization, monic projections, coordinate translations, and every collected quartic coefficient. It also independently replays every retained natural projection, including rational sum-of-squares nonnegativity certificates. Deliberate corruptions must be rejected with assertions disabled as well as enabled.

The architecture checks cover six fully rebuilt variants, finite factorial controls, lower-digit coefficient bounds, and pure-radix regressions. The unbounded fixed-point, carry, CRT, and Horn-flow proofs are given above; finite tests do not supply those proofs. The incoming comparison checks reconstruct ten supplied quartics and four eleven-variable circuit topologies from literal data without importing incoming programs.

\subsection{Manifest and source discipline}
The top-level manifest binds every delivered file other than itself. It records byte sizes and SHA-256 values. Separate source-binding receipts preserve the imported executable closure, actual runtime data, comparison snapshot, original coefficient provenance, and license notices. The archive deliberately excludes original bulk ZIPs, experimental endpoints, temporary replay outputs, generated caches, and unused historical equations. The dominated $446$-integer companion mentioned in the complete signed proof is an audited comparison, not a delivered endpoint; its referenced equation paths are intentionally absent. The six natural architecture variants and all data needed by the principal $442$-integer compiler are retained. Retained incoming receipts describe the reviewed acquisition, not a claim that their complete original archives are embedded.

Every principal scalar coefficient is supplied as gzip-compressed minimal unsigned little-endian bytes. The tuple binding specifies the already-shifted convention for $M_F$; a complete printed-seed recipe, machine transition table, and tile data reconstruct it. Fixed leaves therefore do not represent hidden existential variables or a request to fetch data later.

Verification results, exact endpoint hashes, manifest counts, command counts, and byte-equality counts are recorded in the accompanying release receipt. Passing them certifies the stated finite artifacts. It does not establish the historical native/controller calibration afresh, produce a new internal ZFC derivation, or decide whether any displayed equation has a solution.

\section{Further research and shortcuts that fail}
The improvements leave several concrete research directions. None is asserted as a theorem or an optimality result here.
\begin{enumerate}
\item Seek a lower-degree carrier with the same no-rational-root and every-modulus-root properties while preserving the quantitative powered-projection bound. Reducing the formal carrier degree alone is insufficient if it increases the strong-scale or norm degree.
\item Investigate whether additional common prime factors in the auxiliary modulus can be absorbed into a supported-modulus interpolation family. The strong pair must remain fixed after a residue is prescribed, and the full squared-modulus period must be proved.
\item Search the Pareto frontier of selective even-index choices and residue carriers. Doubling all indices is not a harmless uniform simplification: the retained ordinary $t$ seed is what makes the one-step $t+1$ addition valid and helps bootstrap the radix.
\item Optimize literal monic projection and signed factorization jointly, while requiring full coefficient reconstruction and explicit inverse substitutions. A smaller residual count alone does not guarantee fewer variables, terms, or a lower height.
\item Formalize the new unbounded interfaces in a proof assistant, keeping the already named native and C222 calibration obligations visible. This would be a separate deliverable from the exact finite algebra checks supplied here.
\end{enumerate}

Several tempting shortcuts are invalid. A square product does not imply that both factors are squares without coprimality and positivity. A powered congruence $H\mid K^5$ does not imply $H\mid K$; the strict size argument is essential. A rational ordinate of an adjacent Pell equation cannot simply be declared integral; descent is required. The original period $4m_0$ does not generally suffice modulo $f^2$; the corrected factor $f$ must be retained. The square-output factorial value cannot be assumed before proving radix carry bounds and the ceiling fixed point. Replacing a remainder by $1+qu$ is not a bijective change of arbitrary integer coordinates. Residue one cannot orient a signed power when $q=1$ or $2$. Finally, changing a seed index does not preserve its old private Pell witnesses; existence must be proved by reconstructing them.

The supplied family already includes a quartic in $37$ unrestricted integers with enormous coefficients, and ten-integer equations using a cited Sun reduction. Accordingly, neither $442$ variables for the direct quartic nor eleven variables for the norm is advertised as a global variable-count record. The purpose is a transparent, fully specified and reproducible improvement in the stated multidimensional objectives.

\end{document}
